%% 第13回編集委員会（物理第1巻）：組 P1-2
%% p1-03.tex  力-3　ぶつかる・噴き出す：運動量と衝突（★4）
%% 数値はすべて Python で計算した：scratchpad/phys1/力-3/calc3.py（g＝9.8 m/s²）
%% 出典（本文で使った実機の値・事実）
%%   ・バットとボールがふれている時間 約0.7 ms：D. Russell, Physics of Baseball & Softball Bats, Penn State
%%       https://www.acs.psu.edu/drussell/bats/impulse.html （"The contact time between bat and ball is about 0.7 milliseconds"）
%%   ・テニスボールのはね返り：254 cm から落として 134.62 cm より高く 147.32 cm より低い（ITF の規則 Rule 3 / Appendix I）
%%       https://tennisone.tennisplayer.net/content/rules/rule3.html → e＝0.73〜0.76
%%   ・ゴルフのクラブとボールの反発係数の上限 0.830：USGA "Notice: Spring-Like Effect" (2002-08-06)
%%       https://digital-pd.usopen.com/content/dam/usga/pdf/Equipment/Archived-Notices/8-6-02%20Spring-Like%20Effect.pdf
%%   ・プロ野球の公式球の反発係数の目標値 0.4134：日本野球機構 2015年2月3日発表 https://npb.jp/npb/20150203release.html
%%   ・JNCAP のフルラップ前面衝突試験は時速 55 km：国土交通省 自動車アセスメント
%%       https://www.mlit.go.jp/jidosha/anzen/02assessment/car_h25/test_crash.html
%%   ・LE-7A（H-IIA・H-IIB の1段エンジン）の比推力 440 秒（真空）：三菱重工業 製品ページ
%%       https://www.mhi.com/business/products-services/space-defense/rocket-engines-testing/h2a-h2b-first-stage-engine-le-7a
%%       → 噴き出す速さ u＝440 s×9.81 m/s²≒4.3 km/s（計画書の 4.4 km/s を 4.3 km/s に直した）
%%   ・ロケットの式：ツィオルコフスキー（1903年の論文）。それより前にイギリスの W. ムーアも導いていた（1810年代）
%%   ・ニューヨーク・タイムズ紙の1920年1月13日の社説（ゴダードへの誤解）と、1969年7月17日の訂正記事
%%   ・密な粒の流れの式（μ(I) レオロジー）：Jop, Forterre, Pouliquen, Nature 441, 727 (2006)
\makeatletter
% ベクトル：本文は矢印（第13回の追加の決まり）。太字 \symbf は「専門書をのぞいてみよう」だけ
\newcommand\pxdv[1]{\vec{#1}}
% 台車（横から）  \pxdcart{(左下)}{はば mm}{色}{文字}
\newcommand\pxdcart[4]{\begin{scope}[shift={#1}]
  \filldraw[fill=#3!30, draw=#3!70!black, line width=0.7pt, rounded corners=0.8mm] (0,2.2) rectangle (#2,9.5);
  \fill[bkInk!80] (2.6,1.6) circle (1.6) ({#2-2.6},1.6) circle (1.6);
  \fill[bkLine] (2.6,1.6) circle (0.6) ({#2-2.6},1.6) circle (0.6);
  \node[font=\sffamily\footnotesize, text=bkInk] at ({#2/2},5.85) {#4};\end{scope}}
% ボール  \pxdball{(中心)}{半径}{色}
\newcommand\pxdball[3]{\filldraw[fill=#3!35, draw=#3!75!black, line width=0.7pt] #1 circle (#2);}
% ロケット（右向き）  \pxdrocket{(左のはし)}{長さ}
\newcommand\pxdrocket[2]{\begin{scope}[shift={#1}]
  \filldraw[fill=white, draw=bkInk, line width=0.8pt] (0,-3.2) -- ({#2-7},-3.2) .. controls ({#2-2},-3.2) and (#2,-1) .. (#2,0)
     .. controls (#2,1) and ({#2-2},3.2) .. ({#2-7},3.2) -- (0,3.2) -- cycle;
  \filldraw[fill=bkRed!45, draw=bkRed!70!black, line width=0.6pt] (0,3.2) -- (4,3.2) -- (0,6.5) -- cycle;
  \filldraw[fill=bkRed!45, draw=bkRed!70!black, line width=0.6pt] (0,-3.2) -- (4,-3.2) -- (0,-6.5) -- cycle;
  \fill[bkBlue!45] ({#2*0.62},0) circle (1.4);\end{scope}}
% 噴き出すガス  \pxdgas{(右のはし)}{長さ}
\newcommand\pxdgas[2]{\begin{scope}[shift={#1}]
  \foreach \yy/\xx in {0/0, 1.6/2.5, -1.6/2, 0.8/6, -0.9/7.5, 0/10.5, 1.4/12, -1.5/13} {
    \ifdim\xx pt<#2pt \fill[bkOrange!55] (-\xx-1.6,\yy) circle (1.3);\fi}\end{scope}}
\makeatother

\phchapter{3}{ぶつかる・噴き出す：\\運動量と衝突}{ロケットは、何もない宇宙でどうやって進む？}{\bklevel{4}}{ぶつかる・噴き出す：運動量と衝突}
\bkmichishirube{4-11,5-12,5-14,8-6,力-2}{力-6,力-13,熱-4}

\begin{goal}
\begin{itemize}[label={\color{bkTheme}●}, leftmargin=1.6\zw, itemsep=0pt, topsep=1pt]
\item \bklevel{4}\enspace 力積と運動量から、ぶつかるときの平均の力を求められる。
\item \bklevel{4}\enspace 運動量保存と反発係数から、衝突のあとの速さを求められる。
\item \bklevel{4}\enspace 質量が変わるロケットの速さを求められる。
\end{itemize}
\end{goal}

\phkoushiki{%
\phk{作用・反作用の法則（運動量保存と同じこと）}{中学理科}{ph:力-3:sayou}
\phk{$p=mv$、$F\Delta t=mv'-mv$（運動量と力積）}{高校物理}{ph:力-3:rikiseki}
\phk{重心 $x_{\mathrm G}=\dfrac{m_1x_1+m_2x_2}{m_1+m_2}$}{高校物理}{ph:力-3:juushin}
\phk{運動量保存 $m_1v_1+m_2v_2=m_1v_1'+m_2v_2'$}{高校物理}{ph:力-3:hozon}
\phk{反発係数 $e$、衝突で失われる運動エネルギー}{高校物理}{ph:力-3:hanpatsu}
\phk{ななめの衝突、はね返る高さ $h'=e^2h$}{高校物理}{ph:力-3:naname}
\phk{分裂 $m\Delta v=u\Delta m$（ロケット）}{高校物理}{ph:力-3:bunretsu}
}

\begin{hatena}
\noindent\begin{minipage}[c]{0.62\linewidth}\raggedright
ボートは、オールで水を後ろへ押して進みます。人は、地面をけって歩きます。どちらも、何かを押して、その反対に進んでいます。

ところが、宇宙には空気も地面もありません。押す相手が何もないのに、ロケットは、どうやって速くなるのでしょう。
\end{minipage}\hfill
\begin{minipage}[c]{0.35\linewidth}\centering
\begin{tikzpicture}[x=0.8mm, y=0.8mm]
  \fill[bkNavy!85] (-2,-12) rectangle (52,12);
  \foreach \sx/\sy in {3/8, 12/-9, 22/6, 31/-4, 44/9, 48/-8, 8/-2, 38/2} {\fill[white] (\sx,\sy) circle (0.35);}
  \pxdgas{(14,0)}{14}
  \pxdrocket{(14,0)}{26}
  \node[font=\sffamily\bfseries\large, text=white] at (46,0) {?};
\end{tikzpicture}
\end{minipage}
\end{hatena}

\noindent 力-2では、力が物の速度をどう変えるかを、運動方程式で計算しました。この章では、力が「はたらいていた時間」に注目します。ぶつかる・はね返る・噴き出す、のように、ごく短い時間に大きな力がはたらく動きでは、力の大きさを1つ1つ追うよりも、\textbf{運動量}という量を数えるほうが、ずっと楽に答えが出ます。運動量は、力-13 で水の流れが管を押す力を計算するときにも、同じ形で使います。

\begin{omoidasou}
\bkrecall{力-2}{運動方程式 $\dfrac{\dd\pxdv p}{\dd t}=\sum\pxdv F$（$\pxdv p=m\pxdv v$）。作用・反作用：AがBを押す力と、BがAを押す力は、同じ大きさで反対向き。}
\bkrecall{5-12}{力を位置で積分すると仕事。$\displaystyle\int F\,\dd x$ は $F$-$x$ グラフの面積。}
\bkrecall{4-11}{重心は、点の位置の平均。3点なら $\dfrac{\vec a+\vec b+\vec c}{3}$。}
\bkrecall{5-14}{$\displaystyle\int\frac{1}{x}\,\dd x=\ln x+C$（$x>0$）。$\ln$ は自然対数。}
\bkrecall{力-1}{落ちる物の速さ $v^2=2gh$。上へ投げると、上がった高さ $h=\dfrac{v^2}{2g}$。}
\end{omoidasou}

\begin{yattemiyou}[台車どうしで押し合おう]
\textbf{用意するもの}：同じ台車2台、ばね、糸、はさみ、おもり、ものさし\par\vspace{2pt}
\noindent\begin{minipage}[c]{0.56\linewidth}
\begin{enumerate}[label={\bkcircled[bkGreen]{\arabic*}}, leftmargin=2.4\zw, itemsep=1pt, topsep=0pt]
\item 2台のあいだでばねを縮めて、糸でしばる。
\item 糸を切ると、2台は左右に走り出す。スロー動画で撮り、走り出した直後の0.1秒に進んだ距離をはかろう。
\item 片方におもりをのせ、質量を2倍、3倍にして、くり返そう。
\end{enumerate}
\end{minipage}\hfill
\begin{minipage}[c]{0.42\linewidth}\centering
\begin{tikzpicture}[x=0.62mm, y=0.62mm]
  \draw[bkGray, line width=0.8pt] (4,0) -- (88,0);
  \pxdcart{(20,0)}{22}{bkBlue}{A}
  \pxdcart{(50,0)}{22}{bkOrange}{B}
  \draw[bkInk, line width=0.8pt, decorate, decoration={coil, aspect=0.45, segment length=1.2mm, amplitude=1.2mm}] (42.4,5.85) -- (50,5.85);
  \draw[bkRed, line width=0.6pt] (40,9) .. controls (46,11) .. (52,9);
  \draw[bk arrow, bkBlue] (18,5.85) -- (6,5.85);
  \draw[bk arrow, bkOrange] (74,5.85) -- (86,5.85);
  \node[font=\sffamily\footnotesize, anchor=south] at (46,11) {糸};
  \node[font=\sffamily\footnotesize, anchor=north] at (46,-1.5) {ばね};
\end{tikzpicture}
\end{minipage}
\par\vspace{6pt}
ミナさんの結果です（走り出した直後の0.1秒に進んだ距離）。
\begin{bkcenter}
\begin{bktable}{colspec={Q[l,m,wd=34mm] *{3}{Q[c,m,wd=24mm]}}, rowsep=1.5pt, row{1}={font=\sffamily\small}}
Bの質量（Aの何倍） & 1倍 & 2倍 & 3倍 \\
Aが進んだ距離 & \qty{6.0}{cm} & \qty{8.0}{cm} & \qty{9.0}{cm} \\
Bが進んだ距離 & \qty{6.0}{cm} & \qty{4.0}{cm} & \qty{3.0}{cm} \\
\end{bktable}
\bktabcaption{重い台車ほど、ゆっくり走り出す。}
\end{bkcenter}
\end{yattemiyou}

\noindent やってみようの答え：同じ質量なら、同じ距離だけ反対へ進みました。Bを2倍の質量にすると、Bの進んだ距離はAの半分、3倍にすると3分の1です。同じ0.1秒に進んだ距離は速さを表すので、\textbf{質量×速さが、2台で同じ}になっています（$1\times8.0=2\times4.0$、$1\times9.0=3\times3.0$）。この「質量×速さ」が、この章の主役の運動量です。なぜ同じになるのかは、③で導きます。



%% ======================== 物理の道具 ========================
\bksection{doguu}{物理の道具}

\bksubsection{1}{運動量と力積：力×時間で速度が変わる}
\label{ph:力-3:sec1}
\phboxhead{止めるのに、どれだけの力がいる？}
\noindent\begin{minipage}[c]{0.55\linewidth}\raggedright
生卵を\qty{1.25}{m}の高さから落とします。かたい床に落ちると割れますが、ふとんの上に落ちると割れません。どちらも、落ちてきた速さ（約\qty{5}{m/s}）から0まで止まる点は同じです。ちがうのは、\textbf{止まるまでにかかる時間}です。かたい床では一瞬で止まり、ふとんでは、へこみながら、ゆっくり止まります。
\end{minipage}\hfill
\begin{minipage}[c]{0.41\linewidth}\centering
\begin{tikzpicture}[x=1mm, y=1mm]
  \fill[bkGray!45] (0,0) rectangle (22,3);
  \draw[bkInk, line width=0.7pt] (0,3) -- (22,3);
  \filldraw[fill=bkYellow!40, draw=bkOrange!70!black, line width=0.6pt] (11,7.2) ellipse (3 and 4);
  \draw[bk arrow, bkBlue] (11,22) -- (11,12.5);
  \node[font=\sffamily\footnotesize, anchor=north] at (11,-1) {かたい床};
  \filldraw[fill=bkBlue!20, draw=bkBlue!60!black, line width=0.6pt] (30,0) rectangle (54,6);
  \draw[bkBlue!60!black, line width=0.6pt, fill=white] (36,6) .. controls (38,3) and (46,3) .. (48,6);
  \filldraw[fill=bkYellow!40, draw=bkOrange!70!black, line width=0.6pt] (42,7.5) ellipse (3 and 4);
  \draw[bk arrow, bkBlue] (42,22) -- (42,12.8);
  \node[font=\sffamily\footnotesize, anchor=north] at (42,-1) {ふとん};
\end{tikzpicture}
\end{minipage}

\par\vspace{4pt}
力-2の運動方程式を、時間で積分してみましょう。質量 $m$ が変わらないとき、$m\dfrac{\dd v}{\dd t}=\dfrac{\dd(mv)}{\dd t}$ なので、運動方程式は
\pheq{\frac{\dd(m\pxdv v)}{\dd t}=\pxdv F}{kg·m/s$^2$ ＝ N}
と書けます。左の $m\pxdv v$ を、物の\textbf{運動量}（うんどうりょう）といい、$\pxdv p$ で表します。\label{ph:力-3:undouryou}\bkidx{うんどうりょう}{運動量}
\[\pxdv p=m\pxdv v\qquad\text{（運動量＝質量×速度）}\]
運動量は、速度と同じ向きをもつベクトルです。単位は kg·m/s です。一直線上の動きでは、向きを＋と－で表して $p=mv$ と書きます。

\phboxhead{両辺を、はじめの時刻 $t_1$ からあとの時刻 $t_2$ まで積分する}
左辺は「運動量の変化の速さ」なので、積分すると運動量の変化になります（5-11：微分してから積分すると、もとの量の差）。
\pheq{m\pxdv v_2-m\pxdv v_1=\int_{t_1}^{t_2}\pxdv F\,\dd t}{kg·m/s ＝ N·s}
右辺の「力を時間で積分した量」を\textbf{力積}（りきせき）といいます。5-12では、力を\textbf{位置}で積分して仕事を求めました。力積は、力を\textbf{時間}で積分した量です。形は同じで、横軸が位置から時間に変わっただけです。

\begin{bkfigure}
\begin{tikzpicture}
\begin{axis}[bkplot, width=86mm, height=48mm, xmin=0, xmax=12.5, ymin=0, ymax=3.4, xtick={2,10}, xticklabels={$t_1$,$t_2$}, ytick=\empty, xlabel={時間 $t$}, ylabel={力 $F$}]
  \addplot[fill=bkOrange!25, draw=none, domain=2:10, samples=60] {3*sin(deg((x-2)*pi/8))^2} \closedcycle;
  \addplot[bkcurve, domain=2:10, samples=60] {3*sin(deg((x-2)*pi/8))^2};
  \draw[bkGreen!60!black, line width=1.1pt, dashed] (axis cs:2,1.5) -- (axis cs:10,1.5);
  \node[font=\sffamily\small, fill=bkOrange!25, inner sep=1.5pt] at (axis cs:6,0.75) {面積＝力積};
  \node[font=\sffamily\footnotesize, text=bkGreen!50!black, anchor=west] at (axis cs:10.3,1.5) {平均の力 $\overline{F}$};
  \node[font=\sffamily\footnotesize, text=bkBlue!80!black, anchor=south west] at (axis cs:6.6,2.75) {ぶつかっているときの力};
\end{axis}
\end{tikzpicture}
\bkfigcaption{$F$-$t$ 図。ぶつかる力は、0からふえて、また0にもどる。グラフの下の面積が力積で、それが運動量の変化になる。}\label{ph:力-3:fig-ft}
\end{bkfigure}

\begin{phyomi}[$m\pxdv v_2-m\pxdv v_1=\displaystyle\int\pxdv F\,\dd t$]
「あとの運動量－はじめの運動量＝そのあいだに受けた力積」。運動量は、力を受けた「時間の分だけ」変わる。どんなに大きな力でも、はたらく時間が0なら、速度は変わらない。
\end{phyomi}

\phboxhead{力が一定なら：中学・高校の形}
力 $F$ が時間 $\Delta t$ のあいだ一定なら、積分は長方形の面積で、$F\Delta t$ です。一直線上で書くと、
\pheq{F\Delta t=mv'-mv}{N·s ＝ kg·m/s}
\label{ph:力-3:rikiseki}です（$v$ がはじめ、$v'$（ブイ・ダッシュ）があとの速度）。力が変わるときも、$F$-$t$ 図の面積と同じ面積になる長方形の高さを\textbf{平均の力} $\overline{F}$ とすれば、$\overline{F}\Delta t=mv'-mv$ が使えます（図\ref{ph:力-3:fig-ft}の緑の点線）。

\phdim{N·s ＝ kg·m/s$^2$ × s ＝ kg·m/s}\quad 力積と運動量は、同じ単位です。

\begin{phmijika}[ヘルメット・エアバッグ・マット]
止まるまでの運動量の変化 $mv'-mv$ は、同じ物が同じ速さから止まるなら同じです。だから、\textbf{止まるまでの時間 $\Delta t$ を長くすれば、平均の力 $\overline{F}$ は小さくなります}。ヘルメットの中の発泡スチロール、車のエアバッグ、高跳びのマット、ふとんの上の卵は、どれも「ゆっくり止める」ための工夫です。野球のキャッチャーがミットを少し引きながら取るのも同じです。
\end{phmijika}

\begin{phbaisuru}{止まるまでの時間 $\Delta t$}
面積（力積）が同じまま、横はばが2倍になるので、平均の力は半分になります。時間を10倍にすれば、力は10分の1です。
\par\smallskip\noindent\begin{tikzpicture}
\begin{axis}[bkplot, width=100mm, height=50mm, xmin=0, xmax=9, ymin=0, ymax=5, xtick={2,4,8}, ytick={1,2,4}, xlabel={時間（ms）}, ylabel={力（kN）}]
  \addplot[fill=bkRed!20, draw=bkRed, line width=1.4pt, domain=0:2, samples=40] {4*sin(deg(x*pi/2))^2} \closedcycle;
  \addplot[fill=bkBlue!15, draw=bkBlue, line width=1.4pt, domain=0:4, samples=50] {2*sin(deg(x*pi/4))^2} \closedcycle;
  \addplot[fill=bkGreen!15, draw=bkGreen!70!black, line width=1.4pt, domain=0:8, samples=60] {1*sin(deg(x*pi/8))^2} \closedcycle;
  \node[font=\sffamily\footnotesize, text=bkRed!75!black, anchor=west, fill=white, inner sep=1pt] at (axis cs:2.2,4.4) {2 ms：いちばん大きい力 4 kN};
  \node[font=\sffamily\footnotesize, text=bkBlue!80!black, anchor=west, fill=white, inner sep=1pt] at (axis cs:3.9,2.5) {4 ms：2 kN};
  \node[font=\sffamily\footnotesize, text=bkGreen!50!black, anchor=west, fill=white, inner sep=1pt] at (axis cs:6.4,1.5) {8 ms：1 kN};
\end{axis}
\end{tikzpicture}
\par\noindent 3つのグラフの面積（力積）は、どれも同じです。止まるまでの時間が2倍になると、いちばん大きい力も平均の力も半分になります。
\end{phbaisuru}

\begin{phflow}[力積]
\phflowrow{$\dfrac{\dd(m\pxdv v)}{\dd t}=\pxdv F$}{時間で\\積分する}{$m\pxdv v_2-m\pxdv v_1=\displaystyle\int\pxdv F\,\dd t$}
\phflowrow{$m\pxdv v_2-m\pxdv v_1=\displaystyle\int\pxdv F\,\dd t$}{一直線上で\\$F$ が一定}{$F\Delta t=mv'-mv$}
\end{phflow}

%% ---------------------------------------------------------------
\bksubsection{2}{たくさんの物をまとめて見る：質点系と重心}
\label{ph:力-3:sec2}
\phboxhead{2つの物が押し合うと}
\noindent\begin{minipage}[c]{0.5\linewidth}\centering
\begin{tikzpicture}[x=1mm, y=1mm]
  \draw[bkGray, line width=0.8pt] (-2,0) -- (64,0);
  \pxdball{(16,7)}{6}{bkBlue}
  \pxdball{(42,7)}{8}{bkOrange}
  \node[font=\sffamily\footnotesize] at (16,7) {$m_1$};
  \node[font=\sffamily\footnotesize] at (42,7) {$m_2$};
  \draw[bk arrow, bkRed, line width=1.1pt] (21,7) -- (29.5,7);
  \draw[bk arrow, bkRed, line width=1.1pt] (35,11.5) -- (26.5,11.5);
  \node[font=\sffamily\footnotesize, text=bkRed!75!black, anchor=north] at (25,5.3) {$\pxdv F_{21}$};
  \node[font=\sffamily\footnotesize, text=bkRed!75!black, anchor=south] at (31,12.5) {$\pxdv F_{12}$};
  \draw[bk arrow, bkPurple, line width=1.1pt] (16,0.8) -- (16,-9);
  \draw[bk arrow, bkPurple, line width=1.1pt] (42,-1) -- (42,-12);
  \node[font=\sffamily\footnotesize, text=bkPurple!80!black, anchor=east] at (15,-6) {外力 $\pxdv F_1$};
  \node[font=\sffamily\footnotesize, text=bkPurple!80!black, anchor=west] at (43,-8) {外力 $\pxdv F_2$};
\end{tikzpicture}
\end{minipage}\hfill
\begin{minipage}[c]{0.47\linewidth}\raggedright
2つの物（質点1・質点2）を考えます。\par\vspace{2pt}
$\pxdv F_{21}$：1が2から受ける力\\
$\pxdv F_{12}$：2が1から受ける力\par\vspace{2pt}
このように、考えている物どうしがおよぼし合う力を\textbf{内力}、それ以外のもの（地球・床・空気など）から受ける力を\textbf{外力}といいます。
\end{minipage}

\par\vspace{4pt}
1つ1つの物に、運動方程式を書きます。
\[\frac{\dd\pxdv p_1}{\dd t}=\pxdv F_1+\pxdv F_{21},\qquad \frac{\dd\pxdv p_2}{\dd t}=\pxdv F_2+\pxdv F_{12}\]
2つの式を、左辺どうし・右辺どうしでたします。作用・反作用（力-2）から $\pxdv F_{12}=-\pxdv F_{21}$ なので、内力は打ち消し合って消えます。
\pheq{\frac{\dd(\pxdv p_1+\pxdv p_2)}{\dd t}=\pxdv F_1+\pxdv F_2}{N ＝ N}
\textbf{全体の運動量を変えるのは、外力だけ}です。内力は、運動量を一方から他方へ「わたす」だけで、全体の運動量は変えません。物が3つ、100こになっても、内力はすべて2つずつ組になって消えるので、同じ式が成り立ちます。このように、いくつかの物をまとめて考えたものを\textbf{質点系}といいます。

\phboxhead{作用・反作用と運動量保存は、同じこと}
\label{ph:力-3:sayou}反対に読むこともできます。外力がないときに、2つの物の運動量の和がいつも変わらない（$\dfrac{\dd(\pxdv p_1+\pxdv p_2)}{\dd t}=\pxdv 0$）なら、$\dfrac{\dd\pxdv p_1}{\dd t}=-\dfrac{\dd\pxdv p_2}{\dd t}$、つまり $\pxdv F_{21}=-\pxdv F_{12}$ です。\textbf{「作用・反作用の法則」と「運動量の和が変わらないこと」は、同じ内容を2通りに言ったもの}です。

\phboxhead{重心：運動量の和で動く1つの点}
全体の質量を $M=m_1+m_2$ とします。運動量の和を $M$ でわった速度で動く点を考えましょう。
\[\frac{\pxdv p_1+\pxdv p_2}{M}=\frac{m_1\pxdv v_1+m_2\pxdv v_2}{M}=\frac{\dd}{\dd t}\Bigl(\frac{m_1\pxdv r_1+m_2\pxdv r_2}{M}\Bigr)\]
最後の式のかっこの中の位置を、質点系の\textbf{重心}（じゅうしん）といい、$\pxdv r_{\mathrm G}$ と書きます。
\pheq{\pxdv r_{\mathrm G}=\frac{m_1\pxdv r_1+m_2\pxdv r_2}{m_1+m_2}}{m ＝ kg·m/kg}
\label{ph:力-3:juushin}4-11の重心は「位置の平均」でした。ここでは、質量の大きい物ほど重く数える平均（重みつきの平均）です。$x$ 成分だけを書くと、高校で習う形 $x_{\mathrm G}=\dfrac{m_1x_1+m_2x_2}{m_1+m_2}$ になります。重心の速度を $\pxdv v_{\mathrm G}$ とすると、$M\pxdv v_{\mathrm G}=\pxdv p_1+\pxdv p_2$ なので、上の式は
\pheq{M\frac{\dd\pxdv v_{\mathrm G}}{\dd t}=\pxdv F_1+\pxdv F_2}{N ＝ N}
になります。\textbf{重心は、全部の質量が1点に集まって、外力の和を受けているように動きます}。

\begin{phyomi}[$x_{\mathrm G}=\dfrac{m_1x_1+m_2x_2}{m_1+m_2}$]
てこ（シーソー）でつり合う支点の位置。重い物の近くにある。
\end{phyomi}

\begin{phbaisuru}{$m_2$}
$x_1=0$、$x_2=\qty{1}{m}$ に2つの物を置くと、$x_{\mathrm G}=\dfrac{m_2}{m_1+m_2}$（m）です。$m_2$ を $m_1$ の2倍にすると $x_{\mathrm G}=\dfrac{2}{3}$ m、4倍にすると $\dfrac{4}{5}$ m と、重いほうへ近づきます。でも、2つの物のあいだからは出ません。
\par\smallskip\noindent\begin{tikzpicture}
\begin{axis}[bkplot, width=96mm, height=40mm, xmin=0, xmax=10.5, ymin=0, ymax=1.15, xtick={1,2,4,6,8,10}, ytick={0.5,1}, xlabel={$m_2\div m_1$}, ylabel={$x_{\mathrm G}$（m）}]
  \addplot[bkcurve, domain=0:10.2] {x/(1+x)};
  \addplot[bkthin, dashed, domain=0:10.2] {1};
  \fill[bkRed] (axis cs:1,0.5) circle (1.8pt) (axis cs:2,0.6667) circle (1.8pt) (axis cs:4,0.8) circle (1.8pt);
  \node[font=\sffamily\footnotesize, anchor=north west, text=bkRed!75!black] at (axis cs:4.2,0.78) {4倍で 0.8 m};
  \node[font=\sffamily\footnotesize, anchor=south, text=bkGray] at (axis cs:7.5,1.02) {重い物の位置（1 m）};
\end{axis}
\end{tikzpicture}
\end{phbaisuru}

\begin{seikaku}
形のある物（剛体）の重心は、小さなかけらに分けてたし合わせる積分で
\[\pxdv r_{\mathrm G}=\frac{1}{M}\int\pxdv r\,\dd m\]
と書けます。板なら面で、立体なら体積で積分します（8-6 の重積分）。剛体の重心は、力-5 で求めます。
\end{seikaku}

\begin{phflow}[質点系と重心]
\phflowrow{$\dfrac{\dd\pxdv p_1}{\dd t}=\pxdv F_1+\pxdv F_{21}$\\$\dfrac{\dd\pxdv p_2}{\dd t}=\pxdv F_2+\pxdv F_{12}$}{作用・反作用\\$\pxdv F_{12}=-\pxdv F_{21}$}{$M\dfrac{\dd\pxdv v_{\mathrm G}}{\dd t}=\pxdv F_1+\pxdv F_2$}
\phflowrow{$\pxdv r_{\mathrm G}=\dfrac{m_1\pxdv r_1+m_2\pxdv r_2}{m_1+m_2}$}{一直線上の\\$x$ 成分}{$x_{\mathrm G}=\dfrac{m_1x_1+m_2x_2}{m_1+m_2}$}
\end{phflow}

%% ---------------------------------------------------------------
\bksubsection{3}{運動量保存：合体と分裂}
\label{ph:力-3:sec3}
②の式で、外力の和が0なら、右辺が0です。
\pheq{\frac{\dd(\pxdv p_1+\pxdv p_2)}{\dd t}=\pxdv 0\quad\Longrightarrow\quad \pxdv p_1+\pxdv p_2=\text{一定}}{kg·m/s}
外力の和が0のとき、全体の運動量は変わりません。これを\textbf{運動量保存の法則}といいます。\bkidx{うんどうりょうほぞんのほうそく}{運動量保存の法則}一直線上で、ぶつかる前の速度を $v_1$、$v_2$、あとの速度を $v_1'$、$v_2'$ とすると、
\pheq{m_1v_1+m_2v_2=m_1v_1'+m_2v_2'}{kg·m/s}
\label{ph:力-3:hozon}です。速度は、決めた向きを＋、反対を－にして代入します。

\begin{phmodel}[ぶつかる・分かれる一瞬]
\phignore{ぶつかっているあいだの重力・床の摩擦・空気の抵抗の力積。}
\phassume{ぶつかる時間はとても短い。だから、重力や摩擦（外力）の力積 $F\Delta t$ は、ぶつかる力（内力）の力積に比べて無視できる。}
\phlimit{長い時間押し合うとき（すべりながら止まる、など）。外力がとても大きいとき（地面にめりこむくい打ち、など）。}
\end{phmodel}
重力がはたらいていても、ぶつかる一瞬（\qty{0.001}{s}くらい）なら、重力の力積はとても小さいので、ぶつかる直前と直後では運動量保存が使えます。これが、上の札の意味です。

\phboxhead{合体：ぶつかって一体になる}
\noindent\begin{minipage}[c]{0.52\linewidth}\centering
\begin{tikzpicture}[x=0.88mm, y=0.88mm]
  \draw[bkGray, line width=0.8pt] (-2,10) -- (66,10);
  \pxdcart{(2,10)}{18}{bkBlue}{$m_1$}
  \draw[bk arrow, bkBlue] (21,15.85) -- (31,15.85);
  \node[font=\sffamily\footnotesize, anchor=south] at (26,17) {$v_1$};
  \pxdcart{(40,10)}{18}{bkOrange}{$m_2$}
  \node[font=\sffamily\footnotesize, anchor=south] at (49,20.5) {止まっている};
  \node[font=\sffamily\footnotesize, anchor=east, text=bkGray] at (-3,15) {前};
  \draw[bkGray, line width=0.8pt] (-2,-6) -- (66,-6);
  \pxdcart{(14,-6)}{18}{bkBlue}{$m_1$}
  \pxdcart{(32,-6)}{18}{bkOrange}{$m_2$}
  \draw[bk arrow, bkGreen!70!black] (51,-0.15) -- (61,-0.15);
  \node[font=\sffamily\footnotesize, anchor=south] at (56,1) {$V$};
  \node[font=\sffamily\footnotesize, anchor=east, text=bkGray] at (-3,-1) {後};
\end{tikzpicture}
\end{minipage}\hfill
\begin{minipage}[c]{0.45\linewidth}\raggedright
止まっている $m_2$ に、$m_1$ が速度 $v_1$ でぶつかり、連結して一体になりました。運動量保存から
\[m_1v_1=(m_1+m_2)V\]
\[V=\frac{m_1}{m_1+m_2}\,v_1\]
\end{minipage}

\par\vspace{4pt}
\begin{phbaisuru}{止まっている物の質量 $m_2$}
同じ運動量 $m_1v_1$ を、2台分の質量で分け合うので、$m_2$ が大きいほど、一体になったあとの速さ $V$ は小さくなります。$m_2=m_1$ なら $V=\dfrac{1}{2}v_1$、$m_2=2m_1$ なら $\dfrac{1}{3}v_1$ です。
\par\smallskip\noindent\begin{tikzpicture}
\begin{axis}[bkplot, width=96mm, height=40mm, xmin=0, xmax=6.4, ymin=0, ymax=1.15, xtick={1,2,3,4,5,6}, ytick={0.25,0.5,1}, xlabel={$m_2\div m_1$}, ylabel={$V\div v_1$}]
  \addplot[bkcurve, domain=0:6.2] {1/(1+x)};
  \fill[bkRed] (axis cs:1,0.5) circle (1.8pt) (axis cs:2,0.3333) circle (1.8pt) (axis cs:3,0.25) circle (1.8pt);
  \node[font=\sffamily\footnotesize, anchor=south west, text=bkRed!75!black] at (axis cs:1.15,0.52) {同じ質量なら半分};
\end{axis}
\end{tikzpicture}
\end{phbaisuru}

\phboxhead{分裂：止まっていた物が2つに分かれる}
\bkyattemiyouref の台車は、はじめ2台とも止まっていたので、運動量の和は0です。ばねは内力なので、分かれたあとも和は0のままです。
\[m_{\mathrm A}v_{\mathrm A}+m_{\mathrm B}v_{\mathrm B}=0\qquad\text{だから}\qquad \frac{v_{\mathrm A}}{v_{\mathrm B}}=-\frac{m_{\mathrm B}}{m_{\mathrm A}}\]
速さは質量に反比例し、向きは反対です。Bを2倍の質量にすると、Bの速さはAの半分になります。やってみようで、質量×進んだ距離が2台で同じになったわけは、これです。重心は、はじめ止まっていたので、分かれたあとも同じ場所に止まったままです。

%% ---------------------------------------------------------------
\bksubsection{4}{一直線上の衝突：反発係数}
\label{ph:力-3:sec4}
\phboxhead{ぶつかっている一瞬を、ばねで考える}
ぶつかっている一瞬、2つの物は少しへこみ、またもとにもどろうとします。へこんだ所は、2つの物のあいだにはさんだ小さなばねのように働きます。速さ\qty{2}{m/s}の台車が、止まっている同じ質量の台車にぶつかるとき、この小さなばねのモデルで、2台の速度を計算すると、図\ref{ph:力-3:fig-spring}のようになります。

\begin{bkfigure}
\begin{tikzpicture}
\begin{axis}[bkplot, width=100mm, height=54mm, xmin=0, xmax=1.12, ymin=0, ymax=2.75, xtick={0.5,1}, xticklabels={中間,おわり}, ytick={1,2}, ylabel={速度（m/s）}]
  \fill[bkYellow!25] (axis cs:0,0) rectangle (axis cs:0.5,2.35);
  \fill[bkGreen!12] (axis cs:0.5,0) rectangle (axis cs:1,2.35);
  \addplot[bkcurve, domain=0:1, samples=60] {1+cos(deg(x*pi))};
  \addplot[bkcurve2, domain=0:1, samples=60] {1-cos(deg(x*pi))};
  \fill[bkRed] (axis cs:0.5,1) circle (2pt);
  \node[font=\sffamily\footnotesize, anchor=south] at (axis cs:0.25,2.38) {縮む};
  \node[font=\sffamily\footnotesize, anchor=south] at (axis cs:0.75,2.38) {もどる};
  \node[font=\sffamily\footnotesize, text=bkBlue!80!black, anchor=north west, fill=bkYellow!25, inner sep=1pt] at (axis cs:0.03,1.45) {ぶつかった台車};
  \node[font=\sffamily\footnotesize, text=bkOrange!75!black, anchor=south west, fill=bkYellow!25, inner sep=1pt] at (axis cs:0.03,0.62) {ぶつけられた台車};
  \node[font=\sffamily\footnotesize, text=bkRed!75!black, anchor=west, fill=bkGreen!12, inner sep=1pt] at (axis cs:0.6,1.0) {速度がそろう};
  \node[font=\sffamily\footnotesize, anchor=north] at (axis cs:0.5,-0.3) {ぶつかっている時間};
\end{axis}
\end{tikzpicture}
\bkfigcaption{ぶつかっているあいだの2台の速度（同じ質量、ばねで考えたモデル。へこみが完全にもどるとき）。}\label{ph:力-3:fig-spring}
\end{bkfigure}

\begin{enumerate}[label={\bkcircled{\arabic*}}, leftmargin=2.4\zw, itemsep=1pt, topsep=2pt]
\item 縮む段：ぶつかった台車はおそくなり、ぶつけられた台車は速くなる。2台が近づいている（相対速度が正）あいだ、へこみはふえていく。
\item 速度がそろった瞬間（赤い点）：2台は同じ速度で、へこみがいちばん大きい。
\item もどる段：へこみがもとにもどろうとして、2台を引きはなす。
\end{enumerate}
へこみが完全にもとにもどれば、2台は近づいてきたときと同じ速さではなれていきます。でも、本物の物は、へこんだときに熱や音、こわれた所の形の変化に、エネルギーの一部を使ってしまいます。そのため、はなれる速さは、近づく速さより小さくなります。

\phboxhead{反発係数：はなれる速さ÷近づく速さ}
そこで、近づく速さ（ぶつかる前の相対速度の大きさ）と、はなれる速さ（ぶつかったあとの相対速度の大きさ）の比を、
\pheq{e=-\frac{v_1'-v_2'}{v_1-v_2}}{(m/s)÷(m/s)：単位なし}
\label{ph:力-3:hanpatsu}と書き、\textbf{反発係数}（はねかえり係数）といいます。\bkidx{はんぱつけいすう}{反発係数}$-$ がつくのは、ぶつかる前に近づいていた2つが、あとでは反対にはなれていく（相対速度の符号が変わる）からです。$e$ は0から1までの数です。
\begin{itemize}[leftmargin=1.2\zw, itemsep=1pt, topsep=1pt, label={\color{bkTheme}●}]
\item $e=1$：へこみが完全にもどる（\textbf{弾性衝突}）。はなれる速さ＝近づく速さ。
\item $0<e<1$：ふつうの衝突。
\item $e=0$：へこんだまま、くっついて一体になる（\textbf{完全非弾性衝突}）。③の合体です。
\end{itemize}

\begin{phmada}[反発係数は、なぜその値になる？]
反発係数は、ニュートンが実験から見つけた\textbf{経験則}です。運動方程式から導いたものではありません。実際の $e$ は、材料だけでなく、ぶつかる速さ・温度・物の形や大きさによって変わります（たとえば、速くぶつかるほど、金属の球の $e$ は小さくなっていきます）。ぶつかった一瞬に、物の中でどれだけのエネルギーが熱や振動になるかを、材料の性質から理論だけで正確に予測することは、今もむずかしい問題です。そのため、ボールや部品の $e$ は、ふつうは実際にぶつけてはかって決めています。
\end{phmada}

\begin{bkcenter}
\begin{bktable}{colspec={Q[l,m,wd=58mm] Q[c,m,wd=26mm] Q[l,m,wd=36mm]}, rowsep=1.5pt, row{1}={bg=bkTheme!10, font=\sffamily\small}}
ぶつかる物 & $e$ & はかり方・出典 \\
テニスボールとコンクリート & 0.73〜0.76 & 国際テニス連盟の規則の値から計算 \\
ゴルフのクラブとボール & 0.830 以下 & 規則で決めた上限 \\
プロ野球の公式球 & 0.4134（目標） & 日本野球機構が決めた方法 \\
ねんどのかたまりと床 & ほぼ0 & くっついて止まる \\
\end{bktable}
\bktabcaption{反発係数の例。はかり方が決まっている値だけをのせた。}
\end{bkcenter}

\phboxhead{衝突のあとの速度を求める：2つの式を連立する}
運動量保存と反発係数の式は、$v_1'$、$v_2'$ についての連立方程式です（3-6）。
\[m_1v_1'+m_2v_2'=m_1v_1+m_2v_2,\qquad v_1'-v_2'=-e(v_1-v_2)\]
重心の速度 $v_{\mathrm G}=\dfrac{m_1v_1+m_2v_2}{m_1+m_2}$ を使って解くと、次のようになります（2つめの式に $m_2$ をかけて1つめの式とたすと $v_1'$ が、$m_1$ をかけて1つめからひくと $v_2'$ が出ます）。
\[v_1'=v_{\mathrm G}-e\,\frac{m_2}{m_1+m_2}(v_1-v_2),\qquad v_2'=v_{\mathrm G}+e\,\frac{m_1}{m_1+m_2}(v_1-v_2)\]
この式は覚えなくてかまいません。例題では、毎回2つの式を立てて連立します。

\phboxhead{特別な場合}
\begin{itemize}[leftmargin=1.2\zw, itemsep=1pt, topsep=1pt, label={\color{bkTheme}●}]
\item 同じ質量・$e=1$・相手が止まっている（$v_2=0$）：$v_{\mathrm G}=\dfrac{v_1}{2}$ なので、$v_1'=0$、$v_2'=v_1$。\textbf{速度が入れかわります}。ビリヤードで、まっすぐ当てた球がぴたりと止まるのはこのためです。
\item $e=0$：$v_1'=v_2'=v_{\mathrm G}$。合体の式になります。
\item 相手がとても重い（かべ・床、$m_2\gg m_1$、$v_2=0$）：$\dfrac{m_2}{m_1+m_2}\to1$、$v_{\mathrm G}\to0$ なので、$v_1'=-ev_1$。かべではね返る速さは、ぶつかる速さの $e$ 倍です。
\end{itemize}

\phboxhead{衝突で失われる運動エネルギー}
\label{ph:力-3:sonshitsu}中学校で、速さ $v$ で動く質量 $m$ の物は $K=\dfrac{1}{2}mv^2$ の\textbf{運動エネルギー}をもつと習いました（なぜ $\dfrac{1}{2}mv^2$ になるかは、力-4 で導きます）。衝突で運動エネルギーはどれだけ減るのでしょう。

それぞれの速度を、「重心の速度」と「重心から見た速度」に分けます。相対速度を $w=v_1-v_2$ とすると、
\[v_1=v_{\mathrm G}+\frac{m_2}{M}w,\qquad v_2=v_{\mathrm G}-\frac{m_1}{M}w\qquad（M=m_1+m_2）\]
です（2つの式の差が $w$、$m_1$ 倍と $m_2$ 倍の和が $Mv_{\mathrm G}$ になることで確かめられます）。これを運動エネルギーの和に代入して展開すると（3-8）、$v_{\mathrm G}w$ の項は $\dfrac{m_1m_2}{M}v_{\mathrm G}w-\dfrac{m_1m_2}{M}v_{\mathrm G}w=0$ で消えて、
\[\frac{1}{2}m_1v_1^2+\frac{1}{2}m_2v_2^2=\underbrace{\frac{1}{2}Mv_{\mathrm G}^2}_{\text{重心の動き}}+\underbrace{\frac{1}{2}\,\frac{m_1m_2}{m_1+m_2}\,w^2}_{\text{近づき合う動き}}\]
になります。衝突の前後で、重心の速度 $v_{\mathrm G}$ は変わりません（運動量保存）。変わるのは相対速度で、$w$ が $-ew$ になります。だから、失われる運動エネルギーは
\pheq{\Delta K=\frac{1}{2}\,\frac{m_1m_2}{m_1+m_2}\,(1-e^2)(v_1-v_2)^2}{kg·(m/s)$^2$ ＝ J}
です。$e=1$ なら $\Delta K=0$（運動エネルギーが保存する）、$e=0$ のとき失われ方がいちばん大きくなります。でも、$e=0$ でも、重心の動きの分 $\dfrac{1}{2}Mv_{\mathrm G}^2$ は残ります。運動量が保存するので、全部が止まることはできないからです。

\begin{phbaisuru}[0から1まで]{反発係数 $e$}
同じ質量の物が、止まっている物にぶつかるとき、失われる運動エネルギーは、はじめの運動エネルギーの $\dfrac{1-e^2}{2}$ 倍です。$e$ を0から大きくしていくと、$e^2$ で効くので、はじめはゆっくり、1に近づくと急に損失が0に近づきます。
\par\smallskip\noindent\begin{tikzpicture}
\begin{axis}[bkplot, width=96mm, height=42mm, xmin=0, xmax=1.12, ymin=0, ymax=0.6, xtick={0.2,0.4,0.6,0.8,1}, ytick={0.25,0.5}, yticklabels={25\,\%,50\,\%}, xlabel={$e$}, ylabel={失われる割合}]
  \addplot[bkcurve, domain=0:1] {(1-x^2)/2};
  \fill[bkRed] (axis cs:0,0.5) circle (1.8pt) (axis cs:0.5,0.375) circle (1.8pt) (axis cs:0.8,0.18) circle (1.8pt);
  \node[font=\sffamily\footnotesize, anchor=south west, text=bkRed!75!black] at (axis cs:0.82,0.19) {$e=0.8$ で 18\,\%};
  \node[font=\sffamily\footnotesize, anchor=west, text=bkRed!75!black] at (axis cs:0.03,0.53) {合体（$e=0$）で 50\,\%};
\end{axis}
\end{tikzpicture}
\end{phbaisuru}

\begin{phflow}[一直線上の衝突]
\phflowrow{運動量保存＋反発係数}{同じ質量・$e=1$\\相手が止まっている}{速度が入れかわる}
\phflowrow{運動量保存＋反発係数}{相手がとても重い\\（かべ・床）}{$v_1'=-ev_1$}
\phflowrow{重心の動き＋相対の動き}{相対速度が\\$-e$ 倍になる}{$\Delta K=\dfrac{1}{2}\,\dfrac{m_1m_2}{m_1+m_2}$\\$\times(1-e^2)(v_1-v_2)^2$}
\end{phflow}

%% ---------------------------------------------------------------
\bksubsection{5}{ななめの衝突と、弾み続けるボール}
\label{ph:力-3:sec5}
\phboxhead{なめらかな床に、ななめにぶつかる}
\noindent\begin{minipage}[c]{0.5\linewidth}\centering
\begin{tikzpicture}[x=1mm, y=1mm]
  \fill[bkGray!35] (-4,-4) rectangle (60,0);
  \draw[bkInk, line width=0.8pt] (-4,0) -- (60,0);
  \draw[bk arrow, bkBlue] (4,24) -- (27,1);
  \draw[bk arrow, bkGreen!70!black] (29,1) -- (52,12.5);
  \draw[bkBlue!60, dashed, line width=0.6pt] (4,24) -- (27,24) -- (27,1.5);
  \draw[bkGreen!50!black, dashed, line width=0.6pt] (29,12.5) -- (52,12.5) -- (52,1.5);
  \node[font=\sffamily\footnotesize, text=bkBlue!80!black, anchor=south] at (15.5,24.5) {$v_x$};
  \node[font=\sffamily\footnotesize, text=bkBlue!80!black, anchor=east] at (26.5,15) {$v_y$};
  \node[font=\sffamily\footnotesize, text=bkGreen!50!black, anchor=south] at (40.5,13) {$v_x$（同じ）};
  \node[font=\sffamily\footnotesize, text=bkGreen!50!black, anchor=west] at (52.5,7) {$ev_y$};
  \node[font=\sffamily\footnotesize, anchor=north] at (28,-4.5) {なめらかな床};
\end{tikzpicture}
\end{minipage}\hfill
\begin{minipage}[c]{0.47\linewidth}\raggedright
\label{ph:力-3:naname}床がなめらかなら、床がボールを押す力は、床に垂直な向きだけです（摩擦がない）。だから、\textbf{床に平行な向きには力積を受けず、平行な成分 $v_x$ は変わりません}。\par\vspace{2pt}
床に垂直な成分だけが、床との衝突になり、$e$ 倍の速さではね返ります。
\end{minipage}

\par\vspace{4pt}
運動量の式を成分に分けて考えるのが、ここでのこつです（4-11）。$x$ 成分と $y$ 成分は、別々に保存したり、別々に変わったりします。床に $45^\circ$ でぶつかったボール（$v_x=v_y$）が $e=0.5$ ではね返ると、はね返る向きは床から $\tan^{-1}0.5\approx27^\circ$ で、床に近い向きになります。

\phboxhead{弾む高さは $e^2$ 倍}
\label{ph:力-3:takasa}高さ $h$ から落としたボールが床に着く速さは $\sqrt{2gh}$ です（力-1）。床で $e$ 倍の速さ $e\sqrt{2gh}$ ではね返り、上がる高さ $h'$ は、力-1の $h=\dfrac{v^2}{2g}$ から
\pheq{h'=\frac{\bigl(e\sqrt{2gh}\bigr)^2}{2g}=e^2h}{m ＝ m}
です。速さが $e$ 倍なので、高さは $e^2$ 倍です。

\begin{phbaisuru}[いろいろな値]{反発係数 $e$}
$e=0.8$ なら $0.64h$、$e=0.5$ なら $0.25h$ です。$e$ を半分にすると、弾む高さは4分の1になります。$e$ が1より少し小さいだけでも、何回も弾むうちに、高さはどんどん低くなります（下の表計算）。
\end{phbaisuru}

\phboxhead{表計算：弾み続けるボールは、いつ止まる？}
高さ\qty{1.0}{m}から落とした $e=0.8$ のボールの、$n$ 回めの衝突のあとに上がる高さと、空中にいる時間を表にします。落ちはじめてから1回めに床に着くまでの時間は $t_0=\sqrt{\dfrac{2h}{g}}=\sqrt{\dfrac{2\times1.0}{9.8}}\approx\qty{0.452}{s}$ です（力-1）。$n$ 回めにはね返ると、上がる速さは $e^n$ 倍なので、上がって落ちてくるまでの時間は $2t_0e^n$ です。
\begin{bkcenter}
\bksheet{4}{24}{%
1 & $n$（回め） & 高さ（m） & 空中の時間（s） & 着く時刻（s） \\
2 & 0 & 1.000 & 0.452 & 0.452 \\
3 & 1 & 0.640 & 0.723 & 1.175 \\
4 & 2 & 0.410 & 0.578 & 1.753 \\
5 & 3 & 0.262 & 0.463 & 2.215 \\
6 & 4 & 0.168 & 0.370 & 2.585 \\
7 & 5 & 0.107 & 0.296 & 2.882 \\
8 & 6 & 0.069 & 0.237 & 3.118 \\
9 & 7 & 0.044 & 0.189 & 3.308 \\}
\bktabcaption{B3＝B2*0.8\^{}2（高さは $e^2$ 倍）。C3＝2*C2*0.8、C4から下は上の行の0.8倍。D3＝D2＋C3。}
\end{bkcenter}

\begin{bkfigure}
\begin{tikzpicture}
\begin{axis}[bkplot, width=120mm, height=44mm, xmin=0, xmax=4.4, ymin=0, ymax=1.15, xtick={1,2,3,4}, ytick={0.5,1}, xlabel={時刻（s）}, ylabel={高さ（m）}]
  \draw[bkBlue, line width=1.7pt, line width=1.3pt] plot coordinates {(axis cs:0.000,1.000) (axis cs:0.038,0.993) (axis cs:0.075,0.972) (axis cs:0.113,0.938) (axis cs:0.151,0.889) (axis cs:0.188,0.826) (axis cs:0.226,0.750) (axis cs:0.264,0.660) (axis cs:0.301,0.556) (axis cs:0.339,0.438) (axis cs:0.376,0.306) (axis cs:0.414,0.160) (axis cs:0.452,0.000) (axis cs:0.490,0.128) (axis cs:0.528,0.241) (axis cs:0.566,0.340) (axis cs:0.604,0.425) (axis cs:0.642,0.496) (axis cs:0.680,0.553) (axis cs:0.718,0.596) (axis cs:0.756,0.624) (axis cs:0.794,0.638) (axis cs:0.832,0.638) (axis cs:0.870,0.624) (axis cs:0.908,0.596) (axis cs:0.946,0.553) (axis cs:0.984,0.496) (axis cs:1.022,0.425) (axis cs:1.060,0.340) (axis cs:1.098,0.241) (axis cs:1.137,0.128) (axis cs:1.175,0.000) (axis cs:1.213,0.102) (axis cs:1.252,0.189) (axis cs:1.290,0.262) (axis cs:1.329,0.320) (axis cs:1.367,0.364) (axis cs:1.406,0.393) (axis cs:1.444,0.408) (axis cs:1.483,0.408) (axis cs:1.522,0.393) (axis cs:1.560,0.364) (axis cs:1.599,0.320) (axis cs:1.637,0.262) (axis cs:1.676,0.189) (axis cs:1.714,0.102) (axis cs:1.753,0.000) (axis cs:1.791,0.080) (axis cs:1.830,0.146) (axis cs:1.868,0.197) (axis cs:1.907,0.233) (axis cs:1.946,0.255) (axis cs:1.984,0.262) (axis cs:2.023,0.255) (axis cs:2.061,0.233) (axis cs:2.100,0.197) (axis cs:2.138,0.146) (axis cs:2.177,0.080) (axis cs:2.215,0.000) (axis cs:2.257,0.066) (axis cs:2.298,0.116) (axis cs:2.339,0.149) (axis cs:2.380,0.166) (axis cs:2.421,0.166) (axis cs:2.462,0.149) (axis cs:2.503,0.116) (axis cs:2.544,0.066) (axis cs:2.585,0.000) (axis cs:2.628,0.053) (axis cs:2.670,0.088) (axis cs:2.712,0.105) (axis cs:2.755,0.105) (axis cs:2.797,0.088) (axis cs:2.839,0.053) (axis cs:2.882,0.000) (axis cs:2.921,0.038) (axis cs:2.960,0.061) (axis cs:3.000,0.069) (axis cs:3.039,0.061) (axis cs:3.079,0.038) (axis cs:3.118,0.000) (axis cs:3.150,0.024) (axis cs:3.182,0.039) (axis cs:3.213,0.044) (axis cs:3.245,0.039) (axis cs:3.276,0.024) (axis cs:3.308,0.000) (axis cs:3.333,0.016) (axis cs:3.358,0.025) (axis cs:3.384,0.028) (axis cs:3.409,0.025) (axis cs:3.434,0.016) (axis cs:3.459,0.000)};
  \draw[bkRed, line width=1pt, dashed] (axis cs:4.066,0) -- (axis cs:4.066,0.75);
  \node[font=\sffamily\footnotesize, text=bkRed!75!black, anchor=south] at (axis cs:4.066,0.77) {約4.07 s};
\end{axis}
\end{tikzpicture}
\bkfigcaption{弾むたびに、高さは0.64倍、空中の時間は0.8倍になる。弾む回数はかぎりなく多いのに、止まる時刻は決まっている。}\label{ph:力-3:fig-bounce}
\end{bkfigure}

\phboxhead{かぎりなく多い数の和：等比級数}
着く時刻は、ふえかたがだんだん小さくなり、ある値に近づいていきます。全部の時間 $T$ は
\[T=t_0+2t_0e+2t_0e^2+2t_0e^3+\cdots\]
です。$S=1+e+e^2+e^3+\cdots$ とおくと、$T=t_0+2t_0e\,S$ です。$S$ を求めるために、$S$ を $e$ 倍した $eS=e+e^2+e^3+\cdots$ を、$S$ からひきます。うしろの項がすべて消えて、
\[S-eS=1\qquad\text{だから}\qquad S=\frac{1}{1-e}\quad（0\leqq e<1）\]
このように、一定の数（ここでは $e$）をかけていく数をかぎりなくたしたものを、\textbf{無限等比級数}といいます。\bkidx{むげんとうひきゅうすう}{無限等比級数}$e=0.8$ なら $S=5$ です。表計算で実際にたすと、1、1.8、2.44、2.952、…、10項で4.57、20項で4.95と、5に近づいていきます。だから、
\[T=t_0+2t_0e\cdot\frac{1}{1-e}=t_0\,\frac{1+e}{1-e}=0.452\times\frac{1.8}{0.2}\approx\qty{4.07}{s}\]
です。弾む回数はかぎりなく多くても、全部の時間は有限です。

\begin{seikaku}
「かぎりなくたす」と言っても、本当に無限回たすことはできません。正しくは「たす項の数をふやしていったときに、和が近づいていく値」のことです（5-3 の極限と同じ考え）。$S-eS$ の計算も、$n$ 項までの和 $S_n$ で $S_n-eS_n=1-e^n$、$S_n=\dfrac{1-e^n}{1-e}$ とし、$0\leqq e<1$ なら $n$ をふやすと $e^n\to0$ になることから確かめられます。本物のボールでは、弾む高さがとても低くなると、$e$ が一定というモデルは使えなくなります。
\end{seikaku}

\phboxhead{ビリヤード：同じ質量・$e=1$ なら、直角に分かれる}
止まっている球に、同じ質量の球をななめに当てます（$e=1$、摩擦や回転は無視）。運動量保存はベクトルの式で、
\[m\pxdv v=m\pxdv v_1'+m\pxdv v_2'\qquad\text{だから}\qquad \pxdv v=\pxdv v_1'+\pxdv v_2'\]
$e=1$ なので運動エネルギーも保存し、$\dfrac{1}{2}mv^2=\dfrac{1}{2}m{v_1'}^2+\dfrac{1}{2}m{v_2'}^2$、つまり $v^2={v_1'}^2+{v_2'}^2$ です。

\noindent\begin{minipage}[c]{0.44\linewidth}\centering
\begin{tikzpicture}[x=1mm, y=1mm]
  \draw[bk arrow, bkBlue] (0,0) -- (44,0);
  \draw[bk arrow, bkGreen!70!black] (0,0) -- (33,19.05);
  \draw[bk arrow, bkOrange] (33,19.05) -- (44,0.6);
  \draw[bkInk, line width=0.4pt] ($(33,19.05)+(-150:3)$) -- ++(-60:3) -- ++(30:3);
  \node[font=\sffamily\footnotesize, text=bkBlue!80!black, anchor=north] at (22,-1) {$\pxdv v$（当てた球のはじめの速度）};
  \node[font=\sffamily\footnotesize, text=bkGreen!50!black, anchor=south east] at (16,10) {$\pxdv v_1'$};
  \node[font=\sffamily\footnotesize, text=bkOrange!75!black, anchor=west] at (40,11) {$\pxdv v_2'$};
\end{tikzpicture}
\end{minipage}\hfill
\begin{minipage}[c]{0.53\linewidth}\raggedright
左の式は「$\pxdv v_1'$ と $\pxdv v_2'$ をつなぐと $\pxdv v$ になる」、右の式は「その三角形で三平方の定理が成り立つ」ということです。三平方の定理の逆（4-4）から、$\pxdv v_1'$ と $\pxdv v_2'$ のあいだの角は $90^\circ$ です。\par\vspace{2pt}
内積で言うと、$v^2=|\pxdv v_1'+\pxdv v_2'|^2={v_1'}^2+{v_2'}^2+2\pxdv v_1'\cdot\pxdv v_2'$ なので、$\pxdv v_1'\cdot\pxdv v_2'=0$（4-12：直交）。
\end{minipage}

%% ---------------------------------------------------------------
\bksubsection{6}{質量が変わる物：ロケットの式}
\label{ph:力-3:sec6}
\phboxhead{ガスを少しだけ噴き出す}
\noindent\begin{minipage}[c]{0.5\linewidth}\centering
\begin{tikzpicture}[x=1mm, y=1mm]
  \pxdrocket{(16,14)}{30}
  \node[font=\sffamily\footnotesize, anchor=south] at (31,18.5) {質量 $m$};
  \node[font=\sffamily\footnotesize, anchor=east, text=bkGray] at (4,14) {前};
  \pxdrocket{(16,-6)}{30}
  \fill[bkOrange!60] (8,-6) circle (2.2);
  \draw[bk arrow, bkOrange] (6,-6) -- (-4,-6);
  \node[font=\sffamily\footnotesize, anchor=north, text=bkOrange!75!black] at (2,-8) {$\Delta m$、速さ $u$};
  \draw[bk arrow, bkBlue] (48,-6) -- (58,-6);
  \node[font=\sffamily\footnotesize, anchor=south, text=bkBlue!80!black] at (53,-5) {$\Delta v$};
  \node[font=\sffamily\footnotesize, anchor=north] at (31,-11) {$m-\Delta m$};
  \node[font=\sffamily\footnotesize, anchor=east, text=bkGray] at (-5,0) {後};
\end{tikzpicture}
\end{minipage}\hfill
\begin{minipage}[c]{0.47\linewidth}\raggedright
ロケットといっしょに動きながら見ます（はじめ、ロケットは止まって見える）。質量 $m$ のロケットが、質量 $\Delta m$ のガスを、ロケットから見て速さ $u$ で後ろへ押し出し、ロケットは $\Delta v$ だけ速くなります。
\end{minipage}

\par\vspace{4pt}
はじめの運動量は0です。ガスを押し出すときの力は内力なので、運動量保存が使えます。前向きを＋にして、
\[(m-\Delta m)\,\Delta v+\Delta m\,(\Delta v-u)=0\]
です（ガスは、速くなったあとのロケットから見て $u$ で後ろへ出るので、その速度は $\Delta v-u$）。展開して整理すると、
\pheq{m\,\Delta v=u\,\Delta m}{kg·m/s ＝ m/s·kg}
\label{ph:力-3:bunretsu}になります。高校で習う「分裂」の式です。押し出したガスの運動量 $u\Delta m$ の分だけ、ロケットは反対向きの運動量を受けとります。

\phboxhead{\bkhatenaref の答え}
ロケットは、空気や地面を押しているのではありません。\textbf{自分が積んできた燃料を燃やしてできたガスを、後ろへ押し出しています}。ロケットがガスを後ろへ押す力の反作用として、ガスがロケットを前へ押します。ロケットとガスを合わせた運動量は、はじめと同じ0のままです。押す相手を自分で持っていくので、何もない宇宙でも進めます。むしろ、まわりの空気がないほうが、空気の抵抗がない分だけ有利です。

\phboxhead{何回もくり返すと：プログラムで足してみる}
燃料を一度に全部ではなく、小分けにして何回も噴き出すと、どうなるでしょう。はじめ\qty{10}{t}（そのうち燃料\qty{9}{t}）のロケットが、$u=\qty{3.0}{km/s}$ で燃料を噴き出すとします。燃料を $N$ 回に等分して、1回ごとに $\Delta v=u\dfrac{\Delta m}{m}$ を足していきます（$m$ はその回に噴き出す直前の質量）。
\begin{phcode}
u = 3.0              \# 噴き出す速さ（km/s）
for N in [1, 2, 10, 100, 1000]:
    m, v = 10.0, 0.0 \# はじめの質量（t）と速さ
    dm = 9.0 / N     \# 1回に噴き出す燃料
    for i in range(N):
        v = v + u * dm / m
        m = m - dm
    print(N, round(v, 3))
\end{phcode}
\begin{phcodeout}
1 2.7
2 3.805
10 5.881
100 6.788
1000 6.896
\end{phcodeout}
小分けにするほど速さは大きくなり、ある値（約\qty{6.91}{km/s}）に近づきます。一度に全部噴き出すと、うしろのほうの燃料も重いロケットを押すのに使われてしまいますが、少しずつ出すと、軽くなったロケットを押すのに使えるからです。本物のロケットは燃料を少しずつ、続けて噴き出すので、この近づく先の値が本当の答えです。

\phboxhead{積分すると：ロケットの式}
「小分けをかぎりなく細かくした先」は、積分で求められます。ロケットの質量が $\dd m$（負の数。ふえた分で表す）だけ変わると、噴き出した燃料は $-\dd m$ なので、$m\,\dd v=-u\,\dd m$、つまり
\[\dd v=-u\,\frac{\dd m}{m}\]
です。はじめの質量 $m_0$ から、燃えつきたときの質量 $m_1$ まで積分します（5-14：$\dfrac{1}{m}$ を積分すると $\ln m$）。
\[\Delta v=-u\int_{m_0}^{m_1}\frac{\dd m}{m}=-u\bigl(\ln m_1-\ln m_0\bigr)\]
\pheq{\Delta v=u\ln\frac{m_0}{m_1}}{m/s ＝ m/s（ln の中は単位なし）}
これを\textbf{ロケットの式}といいます。\bkidx{ろけっとのしき}{ロケットの式}$\dfrac{m_0}{m_1}$ を\textbf{質量比}といいます。上の例では $u=3.0$、$\dfrac{m_0}{m_1}=10$ なので、$\Delta v=3.0\times\ln10=3.0\times2.303\approx\qty{6.91}{km/s}$ で、プログラムの近づく先と同じです。

\begin{phyomi}[$\Delta v=u\ln\dfrac{m_0}{m_1}$]
ロケットがふやせる速さは、「噴き出す速さ」×「質量比の自然対数」。燃料をどれだけの時間で燃やすか（推力の大きさ）にはよらない（重力と空気の抵抗を無視したとき）。
\end{phyomi}

\begin{phbaisuru}{質量比 $\dfrac{m_0}{m_1}$}
質量比を2倍にしても、$\Delta v$ は2倍になりません。$\ln$ のグラフはだんだんねていくので、質量比2で $0.69u$、4で $1.39u$、10で $2.30u$、20で $3.00u$ です。質量比を2倍にするごとに、$\Delta v$ は $u\ln2\approx0.69u$ ずつしかふえません。
\par\smallskip\noindent\begin{tikzpicture}
\begin{axis}[bkplot, width=100mm, height=44mm, xmin=0, xmax=22, ymin=0, ymax=3.5, xtick={2,4,10,20}, ytick={1,2,3}, yticklabels={$u$,$2u$,$3u$}, xlabel={質量比}, ylabel={$\Delta v$}]
  \addplot[bkcurve, domain=1:21, samples=80] {ln(x)};
  \fill[bkRed] (axis cs:2,0.693) circle (1.8pt) (axis cs:4,1.386) circle (1.8pt) (axis cs:10,2.303) circle (1.8pt) (axis cs:20,2.996) circle (1.8pt);
  \node[font=\sffamily\footnotesize, anchor=north west, text=bkRed!75!black] at (axis cs:4.3,1.3) {4 で $1.39u$};
  \node[font=\sffamily\footnotesize, anchor=north west, text=bkRed!75!black] at (axis cs:10.3,2.2) {10 で $2.30u$};
\end{axis}
\end{tikzpicture}
\par\noindent だから、本物のロケットは、燃料を使いきった空のタンクやエンジンを途中で捨てて、軽くしてから次の段に火をつけます（多段ロケット）。
\end{phbaisuru}

\begin{motto}[4]
\phboxhead{2段にすると、なぜ速くなる？（表計算）}
全体\qty{100}{t}（燃料\qty{90}{t}、タンクとエンジン\qty{9}{t}、人工衛星\qty{1}{t}）のロケットを、$u=\qty{3.0}{km/s}$ で比べます。2段では、燃料もタンクも 2:1 に分け、1段めを燃やしきったら、1段めのタンク（\qty{6}{t}）を捨てます。
\begin{bkcenter}
\bksheet{4}{24}{%
1 & & はじめ（t） & 燃えつき（t） & $\Delta v$（km/s） \\
2 & 1段だけ & 100 & 10 & 6.91 \\
3 & 2段の1段め & 100 & 40 & 2.75 \\
4 & 2段の2段め & 34 & 4 & 6.42 \\
5 & 2段の合計 & & & 9.17 \\}
\bktabcaption{E2＝3.0*LN(C2/D2)。2段めのはじめは、40 t から捨てた 6 t をひいた 34 t。}
\end{bkcenter}
燃料の量は同じなのに、2段にすると約\qty{9.2}{km/s}になり、人工衛星の速さ約\qty{7.9}{km/s}をこえます。使い終わった重いタンクを運ばなくてよいからです。
\end{motto}

\phboxhead{推力：1秒あたりに押し出す運動量}
1秒あたりに噴き出すガスの質量を $\dot m$（エム・ドット。kg/s。力-13 でも使う）とすると、1秒あたりにガスへ与える運動量は $\dot mu$ です。運動量の変化の速さは力なので（①）、ロケットを押す力（\textbf{推力}）は
\[F=\dot m\,u\]
です。たとえば、噴き出す速さ\qty{3000}{m/s}で、毎秒\qty{100}{kg}のガスを出すと、
\[F=100\times3000=\qty{3.0e5}{N}\]
です。

\begin{phflow}[ロケット]
\phflowrow{運動量保存\\（ガスとロケット）}{少しだけ\\噴き出す}{$m\Delta v=u\Delta m$}
\phflowrow{$\dd v=-u\dfrac{\dd m}{m}$}{$m_0$ から $m_1$ まで\\積分する}{$\Delta v=u\ln\dfrac{m_0}{m_1}$}
\phflowrow{$\Delta v=u\ln\dfrac{m_0}{m_1}$}{燃料が少ない\\$\Delta m\ll m_0$}{$\Delta v\approx u\dfrac{\Delta m}{m_0}$}
\end{phflow}
\noindent 最後の段の $\approx$ は「ほぼ等しい」です（5-8）。燃料が少しなら、$\ln\dfrac{m_0}{m_0-\Delta m}$ は $\dfrac{\Delta m}{m_0}$ にほぼ等しく（5-17 の近似）、1回だけ噴き出す分裂の式にもどります。

%% ---------------------------------------------------------------
\bksubsection{7}{噴流が板を押す力}
\label{ph:力-3:sec7}
\noindent\begin{minipage}[c]{0.48\linewidth}\centering
\begin{tikzpicture}[x=1mm, y=1mm]
  \fill[bkGray!50] (44,-16) rectangle (48,16);
  \draw[bkInk, line width=0.6pt] (44,-16) rectangle (48,16);
  \fill[bkBlue!25] (0,-2.2) -- (43,-2.2) -- (43,2.2) -- (0,2.2) -- cycle;
  \fill[bkBlue!25] (40.5,2.2) -- (43,2.2) -- (43,14) -- (40.5,14) -- cycle;
  \fill[bkBlue!25] (40.5,-2.2) -- (43,-2.2) -- (43,-14) -- (40.5,-14) -- cycle;
  \draw[bk arrow, bkBlue] (6,0) -- (24,0);
  \draw[bk arrow, bkBlue] (38.5,6) -- (38.5,14);
  \draw[bk arrow, bkBlue] (38.5,-6) -- (38.5,-14);
  \draw[bk arrow, bkRed] (49,0) -- (60,0);
  \node[font=\sffamily\footnotesize, anchor=south] at (14,2.8) {速さ $v$};
  \node[font=\sffamily\footnotesize, anchor=north, text=bkGray] at (14,-3) {断面積 $A$};
  \node[font=\sffamily\footnotesize, anchor=south, text=bkRed!75!black] at (55,1) {$F$};
  \node[font=\sffamily\footnotesize, anchor=east] at (37,11) {横へ};
\end{tikzpicture}
\end{minipage}\hfill
\begin{minipage}[c]{0.49\linewidth}\raggedright
ホースの水を、かべに直角に当てます。水ははね返らずに、かべにそって横へ流れていくとします。\par\vspace{2pt}
1秒あたりにかべに当たる水の体積は、断面積×速さ $Q=Av$（流量。5-12）、質量は $\rho Q$ です（$\rho$ は水の密度、約\qty{1000}{kg/m^3}）。
\end{minipage}

\par\vspace{4pt}
かべに当たった水は、かべに向かう向きの速さ $v$ を失います（横へ流れるので、かべの向きの速度は0）。1秒あたりに失う運動量は $\rho Qv$ です。運動量の変化の速さは力なので、かべは水を $\rho Qv$ の力で押しもどしています。作用・反作用から、水がかべを押す力は
\pheq{F=\rho Qv=\rho Av^2}{kg/m$^3$×m$^3$/s×m/s ＝ N}
です。ロケットの推力 $F=\dot m\,u$ と同じ形で、$\dot m=\rho Q$ です。

\begin{phbaisuru}{水の速さ $v$}
速さを2倍にすると、1秒あたりに当たる水の量も2倍、その1つ1つの運動量も2倍なので、力は4倍になります。$v^2$ に比例します。消防のホースが、2人がかりでないとおさえられないほど強く押し返すのは、速い水を大量に出しているからです。
\end{phbaisuru}

\begin{phsaki}{流れの中の運動量の数え方}
ここでは、かべに当たる水だけを考えました。曲がった管の中を流れる水が管を押す力や、ポンプ・水車の羽根が受ける力は、「ある箱（検査体積）に、1秒あたりに入る運動量と出る運動量の差」で計算します。力-13 で、運動量の法則として学びます。ガスの流れが速くなって音の速さに近づくと、気体の密度の変化も考えなければなりません。ロケットのノズルの形のひみつは、熱-4 で学びます。
\end{phsaki}

\begin{phmada}[たくさんの粒がぶつかり合う流れ]
砂時計の砂、サイロから出る小麦、薬の粉のように、たくさんの粒が衝突しながら流れる物（粉体）は、身近なのに、まだよくわかっていない物の1つです。1つ1つの粒の衝突は、この章の運動量保存と反発係数で計算できます。粒の数が数万なら、コンピュータで全部の衝突を追う計算もできます（力-16）。でも、粒がぎっしりつまって流れる場合については、2006年に提案された式が多くの実験に合うことがわかってきたものの、止まっている状態から流れ出す所まで含めて、どんな粉体の流れも1つの式で表す方法は、まだ見つかっておらず、研究が続いています。
\end{phmada}

\begin{phsaki}{光の速さに近い物の運動量}
光の速さに近い粒子では、運動量は $mv$ ではなく、もっと大きくなります。それでも「外力がなければ運動量の和は変わらない」ことは成り立ちます。原子核や素粒子の衝突は、原-11 で学ぶ相対性理論の運動量で計算します。
\end{phsaki}

\begin{senmonsho}[力学の専門書の書き方]
\begin{bkhonbun}
質点系の全運動量を $\symbf P=\sum_i m_i\symbf v_i$ とすると，$\dfrac{\dd\symbf P}{\dd t}=\sum_i\symbf F_i^{(\mathrm e)}$ が成り立つ．内力の和は作用・反作用の法則により $\symbf 0$ となる．外力の和が $\symbf 0$ ならば $\symbf P$ は保存される．
\end{bkhonbun}
\bkyaku{太字の $\symbf P$ は、この本の矢印 $\vec P$ と同じベクトル（4-13）。$\sum_i$ は「$i=1,2,3,\cdots$ の全部についてたす」（5-10 のシグマ）。$\symbf F_i^{(\mathrm e)}$ は「$i$ 番めの物が受ける外力」（e は external＝外の）。この章の②と③を、物がいくつあっても使える形で書いたもの。}
\end{senmonsho}

\begin{taisetsu}
\def\tsl#1{\makebox[6.4\zw][l]{\sffamily\bfseries\small\color{\bkTT}#1}}
\tsl{運動量}$\pxdv p=m\pxdv v$（kg·m/s）。向きをもつ\\[3pt]
\tsl{力積}$\displaystyle\int\pxdv F\,\dd t$＝運動量の変化。一定なら $F\Delta t=mv'-mv$\\[3pt]
\tsl{重心}$x_{\mathrm G}=\dfrac{m_1x_1+m_2x_2}{m_1+m_2}$。外力の和で動く\\[3pt]
\tsl{運動量保存}外力の和が0なら $m_1v_1+m_2v_2=m_1v_1'+m_2v_2'$\\[3pt]
\tsl{反発係数}$e=-\dfrac{v_1'-v_2'}{v_1-v_2}$。かべ・床では $v'=-ev$、$h'=e^2h$\\[3pt]
\tsl{ロケット}$\Delta v=u\ln\dfrac{m_0}{m_1}$。推力 $F=\dot m\,u$。噴流 $F=\rho Qv$
\end{taisetsu}

\begin{kotoba}
\small
\bkword{運動量}{質量×速度。ぶつかり合っても、外力がなければ和が変わらない量。}
\bkword{力積}{力を時間で積分した量。$F$-$t$ 図の面積。運動量の変化に等しい。}
\bkword{内力・外力}{考えている物どうしの力と、それ以外から受ける力。}
\bkword{重心}{質量で重みをつけた位置の平均。外力の和を受けて動く点。}
\bkword{反発係数}{はなれる速さ÷近づく速さ。実験で決める数（経験則）。}
\bkword{質量比}{ロケットの、はじめの質量÷燃えつきたときの質量。}
\end{kotoba}

\bkkigocards{\bkkigocard{$\pxdv p$}{ピー}{運動量 $m\pxdv v$。圧力の $p$（力-12）とは別の量}{momentum}{力-3}}
{\bkkigocard{$e$}{イー}{反発係数（0〜1）。自然対数の底 $e$（5-14）とは別}{coefficient of restitution}{力-3}}

%% ======================== 例題 ========================
\Needspace{10\baselineskip}\bksection{reidai}{例題でたしかめよう}
\label{ph:力-3:ex1}
\begin{phreidai}[バットでボールを打ち返す]
質量\qty{0.145}{kg}の野球のボールが、速さ\qty{40}{m/s}で飛んできました。バットで打ち返すと、ボールは反対向きに速さ\qty{50}{m/s}で飛んでいきました。バットとボールがふれていた時間は\qty{0.70}{ms}（\qty{0.00070}{s}）です。バットがボールを押した平均の力を求めましょう。
\phbamen
ボールの速度が、ごく短い時間で反対向きに変わります。ふれている間の力は変わり続けるので、力の大きさを1つ1つ追うことはできません。わかっているのは、はじめとあとの速度と、時間です。
\phmushi
ボールは点と考え、回転は無視します。ふれている間にボールにはたらく重力も無視します（あとで確かめます）。動きは一直線上とします。
\phzu
\noindent\begin{minipage}[c]{0.5\linewidth}\centering
\begin{tikzpicture}[x=1mm, y=1mm]
  \pxdball{(30,0)}{2.6}{bkGray}
  \draw[bk arrow, bkBlue] (52,4) -- (34,4);
  \node[font=\sffamily\footnotesize, anchor=south] at (43,5) {前 $v=-40$};
  \draw[bk arrow, bkGreen!70!black] (34,-4) -- (60,-4);
  \node[font=\sffamily\footnotesize, anchor=north] at (47,-5) {後 $v'=+50$};
  \draw[bkOrange!70!black, line width=3pt, line cap=round] (25,-9) -- (27,10);
  \draw[bk arrow, bkInk, line width=0.9pt] (2,-12) -- (14,-12);
  \node[font=\sffamily\footnotesize, anchor=west] at (15,-12) {＋の向き};
\end{tikzpicture}
\end{minipage}\hfill
\begin{minipage}[c]{0.47\linewidth}\raggedright
打ち返した向き（投手のほう）を＋にします。はじめの速度 $v=\qty{-40}{m/s}$、あとの速度 $v'=\qty{+50}{m/s}$、時間 $\Delta t=\qty{0.00070}{s}$、平均の力を $\overline{F}$ とします。
\end{minipage}
\phriyu
力が変わり続けても、平均の力と時間をかけた力積は、運動量の変化に等しい（①）。時間がわかっているので、この式が使えます。
\phtateru
\[\overline{F}\Delta t=mv'-mv\]
\phtoku
運動量の変化は、
\[mv'-mv=0.145\times50-0.145\times(-40)\]
\[=0.145\times90=\qty{13.05}{kg.m/s}\]
\[\overline{F}=\frac{13.05}{0.00070}\approx\qty{1.9e4}{N}\]
\phtashikame
単位：kg·m/s÷s＝kg·m/s$^2$＝N で、力の単位です。\\
大きさ：ボールの重さ（重力）$0.145\times9.8\approx\qty{1.4}{N}$ の約1万3000倍です。とても大きな力ですが、ほんの一瞬なので、ボールは大きくへこみながらはね返ります。\\
無視したことの確かめ：ふれている間の重力の力積は、
\[1.4\times0.00070\approx\qty{0.001}{N.s}\]
で、\qty{13}{N.s}に比べて無視できます。\\
極端な場合：バットでボールを止めるだけ（$v'=0$）なら、運動量の変化は $0.145\times40=5.8$ で、力は半分より小さくなります。打ち返すと、止める分と、反対向きに飛ばす分の両方がいるのです。
\kotae 約\qty{1.9e4}{N}（約\qty{19}{kN}）
\end{phreidai}
\begin{jibun}
質量\qty{0.057}{kg}のテニスボールが、速さ\qty{30}{m/s}で飛んできました。ラケットで、反対向きに速さ\qty{40}{m/s}で打ち返しました。ふれていた時間を\qty{5.0}{ms}として、平均の力を求めましょう。
\bkjans{$0.057\times(40+30)=3.99$、$3.99\div0.0050\approx\qty{8.0e2}{N}$（約\qty{800}{N}）}
\end{jibun}

\label{ph:力-3:ex2}
\begin{phreidai}[ヘルメットは、なぜ頭を守る？]
頭（質量\qty{5.0}{kg}とする）が、速さ\qty{10}{m/s}で地面にぶつかって止まります。ヘルメットがないと\qty{0.010}{s}で止まり、ヘルメットの中のやわらかい材料がつぶれると\qty{0.050}{s}で止まるとします。頭が受ける平均の力を、それぞれ求めましょう。
\phbamen
同じ頭が、同じ速さから0まで止まります。止まるまでの時間だけがちがいます。
\phmushi
頭と首のつながりは考えず、頭を\qty{5.0}{kg}の1つの物とします。ぶつかる間の重力は無視します。
\phzu
ぶつかる向きを＋にします。$m=\qty{5.0}{kg}$、$v=\qty{10}{m/s}$、$v'=0$、止まるまでの時間 $\Delta t$。
\phriyu
止まるまでの運動量の変化は、どちらも同じです。平均の力は、力積÷時間で求められます（①の2倍にすると？の考え）。
\phtateru
\[\overline{F}\Delta t=mv'-mv=0-5.0\times10=\qty{-50}{kg.m/s}\]
\phtoku
\[\Delta t=\qty{0.010}{s}\text{ のとき}\quad \overline{F}=\frac{-50}{0.010}=\qty{-5.0e3}{N}\]
\[\Delta t=\qty{0.050}{s}\text{ のとき}\quad \overline{F}=\frac{-50}{0.050}=\qty{-1.0e3}{N}\]
－は、進む向きと反対向きの力（止める向き）という意味です。
\phtashikame
大きさ：時間が5倍になると、力は5分の1です。加速度にすると、$\dfrac{10}{0.010}=\qty{1000}{m/s^2}$（$g$ の約100倍）と $\qty{200}{m/s^2}$（約20倍）です。\\
極端な場合：止まる時間を0に近づけると、力はかぎりなく大きくなります。かたい物どうしがぶつかるとあぶないのは、このためです。
\kotae ヘルメットなし約\qty{5.0}{kN}、ヘルメットあり約\qty{1.0}{kN}
\end{phreidai}
\begin{jibun}
質量\qty{0.060}{kg}の卵を、高さ\qty{1.25}{m}から落とします。床に着く直前の速さを $v=\sqrt{2gh}$ で求め、床で\qty{0.0020}{s}で止まるときと、ふとんで\qty{0.050}{s}で止まるときの、平均の力の大きさを求めましょう。
\bkjans{$v=\sqrt{2\times9.8\times1.25}\approx\qty{4.9}{m/s}$。床：$0.060\times4.95\div0.0020\approx\qty{1.5e2}{N}$、ふとん：約\qty{5.9}{N}}
\end{jibun}

\label{ph:力-3:ex3}
\begin{phreidai}[砂場とマットに着地する]
質量\qty{50}{kg}の人が、速さ\qty{4.0}{m/s}で真下に着地して止まります。止まるまでの時間は、砂場で\qty{0.10}{s}、厚いマットで\qty{0.40}{s}です。足が受ける力の $F$-$t$ 図を三角形とみて、いちばん大きい力を、それぞれ求めましょう。
\phbamen
例題2と同じく止まるまでの時間がちがいますが、こんどは時間が長いので、止まる間の重力も考えます。足が受ける力の形（$F$-$t$ 図）も決めて、いちばん大きい力を求めます。
\phmushi
体を1つの物（質点）とし、ひざの曲げ方などは考えません。力の形を、0から直線的にふえて、また0にもどる三角形とします。
\phzu
\noindent\begin{minipage}[c]{0.5\linewidth}\centering
\begin{tikzpicture}
\begin{axis}[bkplot, width=58mm, height=38mm, xmin=0, xmax=0.47, ymin=0, ymax=5.6, xtick={0.1,0.4}, ytick={2,5}, xlabel={$t$（s）}, ylabel={$N$（kN）}]
  \filldraw[fill=bkRed!20, draw=bkRed, line width=1.3pt] (axis cs:0,0) -- (axis cs:0.05,4.98) -- (axis cs:0.1,0) -- cycle;
  \filldraw[fill=bkBlue!15, draw=bkBlue, line width=1.3pt] (axis cs:0,0) -- (axis cs:0.2,1.98) -- (axis cs:0.4,0) -- cycle;
  \node[font=\sffamily\footnotesize, text=bkRed!75!black, anchor=west] at (axis cs:0.075,4.6) {砂場};
  \node[font=\sffamily\footnotesize, text=bkBlue!80!black, anchor=south] at (axis cs:0.27,2.0) {マット};
\end{axis}
\end{tikzpicture}
\end{minipage}\hfill
\begin{minipage}[c]{0.47\linewidth}\raggedright
上向きを＋にします。床（砂・マット）が足を押す力を $N$、いちばん大きい力を $N_{\max}$、止まるまでの時間を $\Delta t$ とします。体には、上向きの $N$ と、下向きの重力 $mg$ がはたらきます。
\end{minipage}
\phriyu
力の形がわかれば、力積は $F$-$t$ 図の面積（三角形）です。重力も、止まる間ずっとはたらいているので、力積に入れます。
\phtateru
体が受ける力積の和が、運動量の変化です（はじめの速度 $-4.0$、あとの速度0）。
\[\underbrace{\frac{1}{2}N_{\max}\Delta t}_{\text{床の力積}}-mg\Delta t=0-m\times(-4.0)\]
\phtoku
床の力積は $\dfrac{1}{2}N_{\max}\Delta t=m\times4.0+mg\Delta t$ です。
砂場：床の力積は $50\times4.0+50\times9.8\times0.10=200+49=\qty{249}{N.s}$ なので、
\[N_{\max}=\frac{2\times249}{0.10}\approx\qty{5.0e3}{N}\]
マット：床の力積は $200+50\times9.8\times0.40=200+196=\qty{396}{N.s}$ なので、
\[N_{\max}=\frac{2\times396}{0.40}\approx\qty{2.0e3}{N}\]
\phtashikame
単位：N·s÷s＝N。\\
大きさ：止まる時間は4倍ですが、いちばん大きい力は約2.5分の1です（4分の1ではない）。時間が長いと、そのあいだ体重を支える分（重力の力積）がふえるからです。マットの\qty{2.0}{kN}は、体重 $50\times9.8=\qty{490}{N}$ の約4倍です。\\
極端な場合：$\Delta t$ を短くしていくと、重力の力積は無視できるようになり、例題2と同じ形になります。
\kotae 砂場 約\qty{5.0}{kN}、マット 約\qty{2.0}{kN}
\end{phreidai}
\begin{jibun}
質量\qty{40}{kg}の人が、速さ\qty{3.0}{m/s}で真下に着地し、\qty{0.20}{s}で止まりました。足が受ける力の $F$-$t$ 図を三角形として、重力も入れて、いちばん大きい力を求めましょう。
\bkjans{床の力積 $40\times3.0+40\times9.8\times0.20=198.4$、$N_{\max}=2\times198.4\div0.20\approx\qty{2.0e3}{N}$}
\end{jibun}

\label{ph:力-3:ex4}
\begin{phreidai}[貨車をつなぐ]
質量\qty{20}{t}の貨車が、速さ\qty{1.0}{m/s}で、止まっている質量\qty{30}{t}の貨車にぶつかり、連結しました。連結したあとの速さを求めましょう。また、運動エネルギーはどう変わりますか。
\phbamen
ぶつかって、2台が1つになって動きます（合体）。
\phmushi
レールの摩擦と空気の抵抗は、ぶつかる短い時間では無視します（モデル化の札）。
\phzu
進む向きを＋にします。$m_1=\qty{20}{t}=\qty{2.0e4}{kg}$、$v_1=\qty{1.0}{m/s}$、$m_2=\qty{3.0e4}{kg}$、$v_2=0$、連結後の速さ $V$。（③の合体の図と同じ形）
\phriyu
連結器を押し合う力は内力で、水平な外力は無視できるので、水平の運動量が保存します。
\phtateru
\[m_1v_1+m_2\times0=(m_1+m_2)V\]
\phtoku
\[\num{20000}\times1.0=\num{50000}\times V\qquad V=\qty{0.40}{m/s}\]
運動エネルギーは、
\[\text{前：}\ \frac{1}{2}\times\num{20000}\times1.0^2=\qty{1.0e4}{J}\]
\[\text{後：}\ \frac{1}{2}\times\num{50000}\times0.40^2=\qty{4.0e3}{J}\]
\phtashikame
大きさ：$V$ は $v_1$ の $\dfrac{20}{50}=0.4$ 倍で、③のグラフの「$m_2\div m_1=1.5$」の値と合います。\\
エネルギー：\qty{6.0}{kJ}が、連結器のばねや金属の変形・音・熱になりました。合体は $e=0$ なので、運動エネルギーは保存しません（④）。$\dfrac{1}{2}\times\dfrac{\num{20000}\times\num{30000}}{\num{50000}}\times1.0^2=\qty{6.0e3}{J}$ で、④の式とも合います。
\kotae \qty{0.40}{m/s}。運動エネルギーは\qty{10}{kJ}から\qty{4.0}{kJ}に減る
\end{phreidai}
\begin{jibun}
質量\qty{30}{t}の貨車が速さ\qty{1.2}{m/s}で、止まっている\qty{10}{t}の貨車にぶつかって連結しました。連結後の速さと、減った運動エネルギーを求めましょう。
\bkjans{$\num{30000}\times1.2=\num{40000}V$ より \qty{0.90}{m/s}。\qty{21.6}{kJ}から\qty{16.2}{kJ}になり、\qty{5.4}{kJ}減る}
\end{jibun}

\label{ph:力-3:ex5}
\begin{phreidai}[スケートボードの上から荷物を投げる]
スケートボードに乗った人（人とボードで\qty{50}{kg}）が、質量\qty{5.0}{kg}の荷物を持って止まっています。荷物を前へ、地面に対して速さ\qty{4.0}{m/s}で投げると、人とボードはどちら向きに、どんな速さで動き出しますか。
\phbamen
止まっていた物が2つに分かれます（分裂）。
\phmushi
ボードの車輪の転がる抵抗と、空気の抵抗を無視します。荷物は水平に投げるとします。
\phzu
\noindent\begin{minipage}[c]{0.5\linewidth}\centering
\begin{tikzpicture}[x=0.85mm, y=0.85mm]
  \draw[bkGray, line width=0.8pt] (-4,0) -- (62,0);
  \fill[bkInk!80] (14,1.2) circle (1.2) (26,1.2) circle (1.2);
  \fill[bkOrange!60!black] (10,2.4) rectangle (30,3.4);
  \bkperson{(20,3.4)}{24}{bkBlue}
  \filldraw[fill=bkYellow!50, draw=bkOrange!70!black, line width=0.6pt] (38,14) rectangle (45,20);
  \draw[bk arrow, bkGreen!70!black] (47,17) -- (59,17);
  \node[font=\sffamily\footnotesize, anchor=south] at (53,18) {4.0 m/s};
  \draw[bk arrow, bkRed] (8,10) -- (-3,10);
  \node[font=\sffamily\footnotesize, anchor=south] at (2,11) {$V$？};
\end{tikzpicture}
\end{minipage}\hfill
\begin{minipage}[c]{0.47\linewidth}\raggedright
荷物を投げる向きを＋にします。人とボード $M=\qty{50}{kg}$、速度 $V$。荷物 $m=\qty{5.0}{kg}$、速度 $v=\qty{4.0}{m/s}$。
\end{minipage}
\phriyu
はじめは全部止まっていて、運動量の和は0です。投げる力は内力なので、和は0のままです（③の分裂）。
\phtateru
\[MV+mv=0\]
\phtoku
\[50V+5.0\times4.0=0\qquad V=-\frac{20}{50}=\qty{-0.40}{m/s}\]
\phtashikame
符号：－なので、投げた向きと反対（後ろ向き）に動きます。\\
大きさ：質量が10倍なので、速さは10分の1です。\\
極端な場合：地面に立っていると動かないのは、人が地面（地球）とつながっていて、地球全体が反対向きの運動量を受けとるからです。地球の質量はとても大きいので、速さはほぼ0です。
\kotae 後ろ向きに\qty{0.40}{m/s}
\end{phreidai}
\begin{jibun}
ボートに乗った人（人とボートで\qty{80}{kg}）が、止まった状態から、質量\qty{2.0}{kg}のボールを前へ、水に対して速さ\qty{10}{m/s}で投げました。ボートはどちら向きに、どんな速さで動き出しますか。水の抵抗は無視します。
\bkjans{$80V+2.0\times10=0$ より $V=\qty{-0.25}{m/s}$。後ろ向きに\qty{0.25}{m/s}}
\end{jibun}

\label{ph:力-3:ex6}
\begin{phreidai}[ボートの上を歩く]
止まっている質量\qty{120}{kg}のボートの上で、質量\qty{60}{kg}の人が、ボートの後ろのはしから前へ\qty{3.0}{m}歩きました。ボートは、水に対してどれだけ動きますか。
\phbamen
人が歩くと、足がボートを後ろへけるので、ボートは後ろへ動きます。人とボートの速さは、歩くあいだ変わり続けるので、速さから距離を求めるのはむずかしそうです。
\phmushi
水の抵抗を無視します。人とボートの動きは一直線上とします。
\phzu
\noindent\begin{minipage}[c]{0.5\linewidth}\centering
\begin{tikzpicture}[x=0.85mm, y=0.85mm]
  \fill[bkBlue!15] (-4,-3) rectangle (64,2);
  \filldraw[fill=bkOrange!30, draw=bkOrange!70!black, line width=0.7pt] (6,2) -- (46,2) -- (50,8) -- (2,8) -- cycle;
  \bkperson{(10,8)}{18}{bkBlue}
  \draw[bk arrow, bkBlue] (12,30) -- (40,30);
  \node[font=\sffamily\footnotesize, anchor=south] at (26,31) {ボートの上を 3.0 m};
  \draw[bk arrow, bkRed] (4,-6) -- (-4,-6);
  \node[font=\sffamily\footnotesize, anchor=west, text=bkRed!75!black] at (5,-6) {ボートは後ろへ $x$};
\end{tikzpicture}
\end{minipage}\hfill
\begin{minipage}[c]{0.47\linewidth}\raggedright
前向きを＋とし、水に対する位置で考えます。ボートが動く距離を $x$（前向きが＋）とすると、人が水に対して動く距離は $3.0+x$ です。
\end{minipage}
\phriyu
水平の外力がないので、②より、重心の速度は変わりません。はじめ止まっているので、\textbf{重心は動きません}。速さを知らなくても、位置だけで式が立ちます。
\phtateru
重心が動かないので、「質量×動いた距離」の和が0です（$x_{\mathrm G}$ の式で、変化の分だけを考える）。
\[60\times(3.0+x)+120\times x=0\]
\phtoku
\[180+60x+120x=0\qquad 180x=-180\qquad x=\qty{-1.0}{m}\]
\phtashikame
符号：ボートは後ろへ\qty{1.0}{m}動き、人は水に対して前へ $3.0-1.0=\qty{2.0}{m}$ 進みました。\\
確かめ：$60\times2.0=120\times1.0$ で、質量×距離が等しく、重心は動いていません。\\
極端な場合：ボートがとても重いと、$x$ は0に近づき、人はほぼ\qty{3.0}{m}進みます。岸に立って歩くのと同じです。
\kotae 後ろへ\qty{1.0}{m}
\end{phreidai}
\begin{jibun}
止まっている質量\qty{150}{kg}のボートの上で、質量\qty{50}{kg}の人が前へ\qty{4.0}{m}歩きました。ボートは水に対してどれだけ動きますか。
\bkjans{$50(4.0+x)+150x=0$ より $x=\qty{-1.0}{m}$。後ろへ\qty{1.0}{m}}
\end{jibun}

\label{ph:力-3:ex7}
\begin{phreidai}[空中で割れる花火]
質量\qty{1.0}{kg}の花火玉が、いちばん高い所（高さ\qty{45}{m}）で水平に速さ\qty{20}{m/s}で動いているときに、2つに割れました。質量\qty{0.40}{kg}のかけらAは、割れた直後に止まり、そのまま真下に落ちました。もう1つのかけらB（\qty{0.60}{kg}）の割れた直後の速さと、2つが地面に落ちる場所を求めましょう。
\phbamen
空中で1つの物が2つに分かれます。割れる力は内力です。割れたあとは、それぞれが放物運動をします。
\phmushi
空気の抵抗と、割れるときに飛び散る小さなかけら・ガスを無視します。
\phzu
水平で、玉の進む向きを＋の $x$ 向きにします。割れた所の真下を $x=0$ とします。Bの割れた直後の速度を $v_{\mathrm B}$（水平）とします。
\phriyu
割れる一瞬は、重力の力積が無視できるので、水平の運動量が保存します（③）。割れたあとは重力だけの放物運動（力-1）で、②より、2つの重心は割れなかったときの玉と同じ放物線をえがきます。
\phtateru
\[1.0\times20=0.40\times0+0.60\times v_{\mathrm B}\]
\phtoku
$v_{\mathrm B}=\dfrac{20}{0.60}\approx\qty{33}{m/s}$。高さ\qty{45}{m}から落ちる時間は、割れた直後のたての速度が0なので、どちらも
\[t=\sqrt{\frac{2\times45}{9.8}}\approx\qty{3.03}{s}\]
です。落ちる場所は、
\[\text{A：}x=0,\qquad \text{B：}x=33.3\times3.03\approx\qty{101}{m}\]
\phtashikame
重心：$\dfrac{0.40\times0+0.60\times101}{1.0}\approx\qty{61}{m}$。割れなかった玉が落ちる場所 $20\times3.03\approx\qty{61}{m}$ と同じです。\\
極端な場合：割れても、かけらの質量が等しく、同じ速さなら、2つとも同じ所に落ちます。
\kotae Bの速さ約\qty{33}{m/s}。Aは真下（$x=0$）、Bは約\qty{1.0e2}{m}先に落ちる
\end{phreidai}
\begin{jibun}
質量\qty{2.0}{kg}の玉が、高さ\qty{20}{m}のいちばん高い所で水平に速さ\qty{10}{m/s}で動いているときに割れ、\qty{1.5}{kg}のかけらは止まって真下に落ちました。残りのかけら（\qty{0.50}{kg}）の割れた直後の速さと、落ちる場所（割れた所の真下からの距離）を求めましょう。
\bkjans{\qty{40}{m/s}。落ちる時間 $\sqrt{40\div9.8}\approx\qty{2.02}{s}$ で、約\qty{81}{m}先}
\end{jibun}

\label{ph:力-3:ex8}
\begin{phreidai}[弾む高さから反発係数を求める]
ボールを高さ\qty{1.0}{m}から落とすと、床ではね返って\qty{0.64}{m}まで上がりました。ボールと床の反発係数を求めましょう。
\phbamen
床はとても重い相手です。高さから、ぶつかる前後の速さがわかります。
\phmushi
空気の抵抗とボールの回転を無視します。床はかたく、動かないとします。
\phzu
床に着く直前の速さ $v$（下向き）、はね返った直後の速さ $v'$（上向き）。落とす高さ $h=\qty{1.0}{m}$、上がった高さ $h'=\qty{0.64}{m}$。
\phriyu
床は動かないので、④の「相手がとても重い」場合で $e=\dfrac{v'}{v}$ です。速さは高さから $\sqrt{2gh}$ で求まります。⑤の $h'=e^2h$ を使えば1行です。
\phtateru
\[h'=e^2h\qquad\text{だから}\qquad e=\sqrt{\frac{h'}{h}}\]
\phtoku
\[e=\sqrt{\frac{0.64}{1.0}}=0.80\]
速さで確かめると、
\[v=\sqrt{2\times9.8\times1.0}\approx\qty{4.43}{m/s}\]
\[v'=\sqrt{2\times9.8\times0.64}\approx\qty{3.54}{m/s}\]
で、$\dfrac{3.54}{4.43}\approx0.80$ です。
\phtashikame
単位：m÷m なので、単位はありません。\\
大きさ：0と1のあいだです。高さの比0.64ではなく、その平方根が $e$ です。\\
極端な場合：$h'=h$ なら $e=1$、まったく弾まなければ $e=0$。
\kotae $e=0.80$
\end{phreidai}
\begin{jibun}
テニスボールを高さ\qty{2.54}{m}からコンクリートの上に落とすと、\qty{1.40}{m}まで上がりました。反発係数を小数第2位まで求めましょう。
\bkjans{$\sqrt{1.40\div2.54}\approx0.74$}
\end{jibun}

\label{ph:力-3:ex9}
\begin{phreidai}[台車どうしの衝突（反発係数0.5）]
なめらかな水平のレールの上で、質量\qty{2.0}{kg}の台車Aが速さ\qty{3.0}{m/s}で、止まっている質量\qty{1.0}{kg}の台車Bにまっすぐぶつかりました。反発係数を0.50として、ぶつかったあとの2台の速度と、失われた運動エネルギーを求めましょう。
\phbamen
一直線上の衝突です。はね返り方は反発係数で決まります。
\phmushi
レールの摩擦と空気の抵抗を無視します。台車は回転しないとします。
\phzu
Aの進む向きを＋にします。$m_1=\qty{2.0}{kg}$、$v_1=\qty{3.0}{m/s}$、$m_2=\qty{1.0}{kg}$、$v_2=0$、あとの速度 $v_1'$、$v_2'$。
\phriyu
知りたい量が2つ（$v_1'$ と $v_2'$）なので、式も2つ要ります。運動量保存と反発係数の式を連立します（④）。
\phtateru
\[2.0\times3.0+1.0\times0=2.0v_1'+1.0v_2'\]
\[v_1'-v_2'=-0.50\times(3.0-0)\]
\phtoku
2つめの式から $v_2'=v_1'+1.5$。1つめに代入して、
\[6.0=2.0v_1'+v_1'+1.5\qquad 3.0v_1'=4.5\]
\[v_1'=\qty{1.5}{m/s},\qquad v_2'=\qty{3.0}{m/s}\]
運動エネルギーは、
\[\text{前：}\frac{1}{2}\times2.0\times3.0^2=\qty{9.0}{J}\]
\[\text{後：}\frac{1}{2}\times2.0\times1.5^2+\frac{1}{2}\times1.0\times3.0^2=\qty{6.75}{J}\]
で、\qty{2.25}{J}失われました。
\phtashikame
運動量：前 6.0、後 $2.0\times1.5+1.0\times3.0=6.0$ で、保存しています。\\
④の式で：$\dfrac{1}{2}\times\dfrac{2.0\times1.0}{3.0}\times(1-0.25)\times3.0^2=\qty{2.25}{J}$ で、同じです。\\
極端な場合：$e=1$ にすると $v_1'=\qty{1.0}{m/s}$、$v_2'=\qty{4.0}{m/s}$ で、失われるエネルギーは0になります。
\kotae A：\qty{1.5}{m/s}、B：\qty{3.0}{m/s}（どちらも同じ向き）。失われた運動エネルギー\qty{2.25}{J}
\end{phreidai}
\begin{jibun}
質量\qty{1.0}{kg}の台車Aが速さ\qty{4.0}{m/s}で、止まっている質量\qty{3.0}{kg}の台車Bにまっすぐぶつかりました。反発係数0.50として、ぶつかったあとの2台の速度と、失われた運動エネルギーを求めましょう。
\bkjans{A：$\qty{-0.50}{m/s}$（はね返る）、B：\qty{1.5}{m/s}。失われた運動エネルギー\qty{4.5}{J}}
\end{jibun}

\label{ph:力-3:ex10}
\begin{phreidai}[車の追突]
質量\qty{1000}{kg}の車が、速さ\qty{15}{m/s}で、止まっている質量\qty{1500}{kg}の車に追突し、2台は一体になって動きました。ぶつかった直後の速さと、失われた運動エネルギーを求めましょう。
\phbamen
自動車の衝突事故の、ぶつかった一瞬だけを考えます。
\phmushi
ぶつかる短い時間のあいだの、タイヤと路面の摩擦を無視します。車の回転も無視します。
\phzu
進む向きを＋。$m_1=\qty{1000}{kg}$、$v_1=\qty{15}{m/s}$、$m_2=\qty{1500}{kg}$、$v_2=0$。一体になったあとの速さ $V$。
\phriyu
一体になるので $e=0$ です。運動量保存（合体）で $V$ を、④の式か、前後の運動エネルギーの差で、失われた分を求めます。
\phtateru
\[1000\times15=(1000+1500)V\]
\phtoku
\[V=\frac{\num{15000}}{2500}=\qty{6.0}{m/s}\]
\[\text{前：}\frac{1}{2}\times1000\times15^2=\qty{1.125e5}{J}\]
\[\text{後：}\frac{1}{2}\times2500\times6.0^2=\qty{4.5e4}{J}\]
失われた運動エネルギーは $\num{112500}-\num{45000}=\qty{6.75e4}{J}$ です。
\phtashikame
④の式：$\dfrac{1}{2}\times\dfrac{1000\times1500}{2500}\times(1-0)\times15^2=\qty{6.75e4}{J}$ で同じ。\\
大きさ：はじめの運動エネルギーの60\,\%が、車体の変形・熱・音になりました。この大きなエネルギーを、車体の前と後ろの部分がつぶれて受け止めます（工学ではこう使う）。
\kotae \qty{6.0}{m/s}、約\qty{68}{kJ}
\end{phreidai}
\begin{jibun}
質量\qty{2000}{kg}のトラックが速さ\qty{10}{m/s}で、止まっている質量\qty{1000}{kg}の車に追突し、一体になりました。直後の速さと、失われた運動エネルギーを求めましょう。
\bkjans{約\qty{6.7}{m/s}。約\qty{33}{kJ}}
\end{jibun}

\label{ph:力-3:ex11}
\begin{phreidai}[ビリヤードの球]
止まっている球に、同じ質量の球を速さ\qty{2.0}{m/s}で当てます（$e=1$）。（1）まっすぐ当てたとき、（2）当てられた球が、はじめの向きから $30^\circ$ の向きへ進んだとき、ぶつかったあとの2つの球の速さを求めましょう。
\phbamen
（1）は一直線上の衝突、（2）は平面の上の衝突です。
\phmushi
台との摩擦と、球の回転を無視します。$e=1$ で、運動エネルギーは保存するとします。
\phzu
（2）は⑤のビリヤードの図と同じです。はじめの速度 $\pxdv v$（大きさ\qty{2.0}{m/s}）、当てた球のあとの速度 $\pxdv v_1'$、当てられた球のあとの速度 $\pxdv v_2'$。
\phriyu
（1）は④の特別な場合（同じ質量・$e=1$）です。（2）は⑤で、運動量保存と運動エネルギー保存から、2つの球は直角に分かれます。
\phtateru
（2）$\pxdv v=\pxdv v_1'+\pxdv v_2'$ で、$\pxdv v_1'\perp\pxdv v_2'$ なので、$\pxdv v$ を斜辺とする直角三角形ができます。
\phtoku
（1）速度が入れかわり、当てた球は止まり、当てられた球は\qty{2.0}{m/s}で進みます。\\
（2）当てられた球の速さは、斜辺×$\cos30^\circ$ で $2.0\times0.866\approx\qty{1.7}{m/s}$。当てた球は、斜辺×$\sin30^\circ$ で $2.0\times0.5=\qty{1.0}{m/s}$、向きははじめの向きから反対側へ $60^\circ$ です。
\phtashikame
エネルギー：$1.73^2+1.0^2=3.0+1.0=4.0=2.0^2$ で、三平方の定理が成り立っています。\\
極端な場合：$30^\circ$ を0に近づけると（1）にもどり、当てた球は止まります。
\kotae （1）当てた球0、当てられた球\qty{2.0}{m/s}　（2）当てられた球 約\qty{1.7}{m/s}、当てた球\qty{1.0}{m/s}
\end{phreidai}
\begin{jibun}
例題の（2）で、当てられた球が、はじめの向きから $60^\circ$ の向きへ進みました。2つの球の速さを求めましょう。
\bkjans{当てられた球 $2.0\times\cos60^\circ=\qty{1.0}{m/s}$、当てた球 $2.0\times\sin60^\circ\approx\qty{1.7}{m/s}$}
\end{jibun}

\label{ph:力-3:ex12}
\begin{phreidai}[弾み続けるボールが止まるまでの時間]
反発係数0.80のボールを、高さ\qty{1.0}{m}からそっと放しました。弾まなくなるまでの時間を求めましょう。
\phbamen
ボールは何回も弾みます。回数はかぎりなく多いのに、全部の時間は決まった値に近づきます（⑤の表計算）。
\phmushi
空気の抵抗と、床にふれている時間を無視します。$e$ は、どの高さでも一定とします。
\phzu
はじめに床に着くまでの時間 $t_0$。$n$ 回めに弾んだあと空中にいる時間 $2t_0e^n$。全部の時間 $T$。
\phriyu
弾むたびに速さが $e$ 倍になるので、空中の時間も $e$ 倍ずつになります。$e$ 倍ずつの数の和なので、⑤の無限等比級数を使います。
\phtateru
\[T=t_0+2t_0e+2t_0e^2+\cdots=t_0+2t_0e\cdot\frac{1}{1-e}=t_0\,\frac{1+e}{1-e}\]
\phtoku
\[t_0=\sqrt{\frac{2\times1.0}{9.8}}\approx\qty{0.452}{s},\qquad T=0.452\times\frac{1.80}{0.20}\approx\qty{4.1}{s}\]
\phtashikame
表計算：⑤の表で、10回めに着く時刻は\qty{3.68}{s}で、\qty{4.1}{s}に近づいていく途中です。\\
極端な場合：$e\to1$ では $1-e\to0$ で、$T$ はかぎりなく長くなります（いつまでも弾む）。$e=0$ では $T=t_0$ で、1回も弾みません。
\kotae 約\qty{4.1}{s}
\end{phreidai}
\begin{jibun}
反発係数0.50のボールを高さ\qty{0.80}{m}から放しました。弾まなくなるまでの時間を求めましょう。
\bkjans{$t_0=\sqrt{2\times0.80\div9.8}\approx\qty{0.404}{s}$、$T=0.404\times\dfrac{1.5}{0.5}\approx\qty{1.2}{s}$}
\end{jibun}

\label{ph:力-3:ex13}
\begin{phreidai}[1段のロケットで宇宙へ行ける？]
燃料を噴き出す速さが\qty{3.0}{km/s}のロケットがあります。質量比（はじめの質量÷燃えつきたときの質量）が10のとき、ふやせる速さを求めましょう。地球をまわる人工衛星の速さ約\qty{7.9}{km/s}（力-2）に届くでしょうか。
\phbamen
燃料を燃やしながら、質量が変わっていく物の速さです。
\phmushi
重力と空気の抵抗を無視します。噴き出す速さは一定とします。
\phzu
噴き出す速さ $u=\qty{3.0}{km/s}$、質量比 $\dfrac{m_0}{m_1}=10$、ふえる速さ $\Delta v$。
\phriyu
質量が変わるので、$m\Delta v=u\Delta m$ を積分したロケットの式（⑥）を使います。
\phtateru
\[\Delta v=u\ln\frac{m_0}{m_1}\]
\phtoku
\[\Delta v=3.0\times\ln10=3.0\times2.303\approx\qty{6.9}{km/s}\]
\phtashikame
大きさ：\qty{7.9}{km/s}に届きません。しかも本物では、重力と空気の抵抗に負ける分があるので、もっと足りなくなります。\\
改良：噴き出す速さを上げると $\Delta v$ は $u$ に比例して大きくなります。液体水素と液体酸素を燃やすエンジン（LE-7A）では $u$ が約\qty{4.3}{km/s}で、同じ質量比10なら $4.3\times2.303\approx\qty{9.9}{km/s}$ です。それでも、空のタンクを途中で捨てる多段ロケットにして、重力や空気の抵抗に負ける分をおぎないます。
\kotae 約\qty{6.9}{km/s}。1段で、この条件では届かない
\end{phreidai}
\begin{jibun}
噴き出す速さ\qty{2.5}{km/s}、質量比8のロケットがふやせる速さを求めましょう（$\ln8\approx2.08$）。
\bkjans{$2.5\times2.08\approx\qty{5.2}{km/s}$}
\end{jibun}

\label{ph:力-3:ex14}
\begin{phreidai}[ホースの水がかべを押す力]
ホースの口から、水が速さ\qty{10}{m/s}で出ています。口の断面積は\qty{2.0}{cm^2}です。この水をかべに直角に当てると、水ははね返らずにかべにそって横へ流れました。水がかべを押す力を求めましょう。水の密度は\qty{1000}{kg/m^3}とします。
\phbamen
次々に当たる水が、かべを押し続けます。
\phmushi
水の広がりや、空気の抵抗、重力で水が曲がることを無視します。
\phzu
⑦の図と同じです。
\[v=\qty{10}{m/s},\quad A=\qty{2.0}{cm^2}=\qty{2.0e-4}{m^2},\quad \rho=\qty{1000}{kg/m^3}\]
とし、かべを押す力を $F$ とします。
\phriyu
1秒あたりに失う運動量が力になります（⑦）。
\phtateru
\[F=\rho Av^2\]
\phtoku
流量 $Q=Av=\num{2.0e-4}\times10=\qty{2.0e-3}{m^3/s}$（1分あたり\qty{120}{L}）。
\[F=1000\times\num{2.0e-3}\times10=\qty{20}{N}\]
\phtashikame
単位：kg/m$^3$×m$^3$/s×m/s＝kg·m/s$^2$＝N。\\
大きさ：\qty{20}{N}は、約\qty{2}{kg}の物の重さくらいです。手でかべの代わりに受けると、しっかり押されるのを感じます。\\
はね返る場合：水がまっすぐはね返ると、運動量の変化は2倍で、力も約2倍になります。
\kotae \qty{20}{N}
\end{phreidai}
\begin{jibun}
口の断面積\qty{1.0}{cm^2}のホースから、水が速さ\qty{15}{m/s}で出ています。かべに直角に当てて横へ流れるとき、かべを押す力を求めましょう。
\bkjans{$1000\times\num{1.0e-4}\times15^2\approx\qty{23}{N}$}
\end{jibun}

\label{ph:力-3:ex15}
\begin{phreidai}[ばねのついた台車の衝突]
なめらかな水平面の上で、前にばね（ばね定数\qty{800}{N/m}、質量は無視）をつけた質量\qty{2.0}{kg}の台車Aが、速さ\qty{3.0}{m/s}で、止まっている質量\qty{1.0}{kg}の台車Bにぶつかります。ばねがいちばん縮んだときの2台の速さと、ばねの縮みを求めましょう。また、ばねがもとにもどって2台がはなれたあとの速度を求めましょう。
\phbamen
④の「ばねで考える衝突」を、本物のばねでやります。ばねが縮むあいだも、もどるあいだも、ばねの力は内力です。
\phmushi
床の摩擦、空気の抵抗、ばねの質量、熱になるエネルギーを無視します。
\phzu
\noindent\begin{minipage}[c]{0.5\linewidth}\centering
\begin{tikzpicture}[x=0.85mm, y=0.85mm]
  \draw[bkGray, line width=0.8pt] (-2,0) -- (64,0);
  \pxdcart{(0,0)}{20}{bkBlue}{A}
  \draw[bkInk, line width=0.8pt, decorate, decoration={coil, aspect=0.45, segment length=1.6mm, amplitude=1.6mm}] (20,5.85) -- (32,5.85);
  \fill[bkInk] (32,3.6) rectangle (33,8.1);
  \pxdcart{(40,0)}{16}{bkOrange}{B}
  \draw[bk arrow, bkBlue] (4,13) -- (16,13);
  \node[font=\sffamily\footnotesize, anchor=south] at (10,14) {3.0 m/s};
\end{tikzpicture}
\end{minipage}\hfill
\begin{minipage}[c]{0.47\linewidth}\raggedright
右向きを＋。$m_1=\qty{2.0}{kg}$、$v_1=\qty{3.0}{m/s}$、$m_2=\qty{1.0}{kg}$、$v_2=0$、$k=\qty{800}{N/m}$。いちばん縮んだときの速さ $V$、縮み $x$。
\end{minipage}
\phriyu
いちばん縮むのは、2台の速度がそろったとき（④の赤い点）です。そのときの速さは合体と同じ式で求まります。減った運動エネルギーは、ばねにたくわえられています。ばねにたくわえられるエネルギーは $\dfrac{1}{2}kx^2$ です（5-12④。力-4 でくわしく）。はなれたあとは、エネルギーが失われないので $e=1$ です。
\phtateru
\[2.0\times3.0=(2.0+1.0)V,\qquad \frac{1}{2}\times2.0\times3.0^2=\frac{1}{2}\times3.0\times V^2+\frac{1}{2}kx^2\]
\phtoku
$V=\qty{2.0}{m/s}$。ばねのエネルギーは $9.0-6.0=\qty{3.0}{J}$ なので、
\[\frac{1}{2}\times800\times x^2=3.0\qquad x^2=0.0075\qquad x\approx\qty{0.087}{m}\]
はなれたあとは、運動量保存と $v_1'-v_2'=-(3.0-0)$ を連立して、$v_1'=\qty{1.0}{m/s}$、$v_2'=\qty{4.0}{m/s}$。
\phtashikame
運動量：$2.0\times1.0+1.0\times4.0=6.0$。エネルギー：$1.0+8.0=\qty{9.0}{J}$ で、もとと同じです。\\
極端な場合：ばねをかたくする（$k$ を大きくする）と縮みは小さくなりますが、速度がそろったときの速さ $V$ は $k$ によりません。
\kotae いちばん縮んだとき2台とも\qty{2.0}{m/s}、縮み約\qty{8.7}{cm}。はなれたあと A：\qty{1.0}{m/s}、B：\qty{4.0}{m/s}
\end{phreidai}
\begin{jibun}
前にばね（\qty{1200}{N/m}）をつけた質量\qty{1.0}{kg}の台車が速さ\qty{4.0}{m/s}で、止まっている\qty{3.0}{kg}の台車にぶつかります。ばねがいちばん縮んだときの速さと、ばねの縮みを求めましょう。
\bkjans{$V=\qty{1.0}{m/s}$。ばねのエネルギー $8.0-2.0=\qty{6.0}{J}$ より縮み\qty{0.10}{m}}
\end{jibun}

\label{ph:力-3:ex16}
\begin{phreidai}[そりに荷物がすべりこむ]
なめらかな氷の上に、質量\qty{20}{kg}のそりが止まっています。質量\qty{10}{kg}の荷物が、速さ\qty{3.0}{m/s}ですべってきて、そりの上にのり、そりの上をすべりながら、やがてそりといっしょに動きました。荷物とそりのあいだの動摩擦係数を0.30として、（1）いっしょに動くときの速さ、（2）そうなるまでの時間、（3）荷物がそりの上をすべった距離を求めましょう。
\phbamen
荷物とそりのあいだには摩擦力がはたらき、しばらくすべります。短い衝突ではありませんが、摩擦力は荷物とそりのあいだの内力です。
\phmushi
氷とそりの摩擦、空気の抵抗を無視します。荷物がのるときの段差は考えません。
\phzu
進む向きを＋。荷物 $m=\qty{10}{kg}$、$v=\qty{3.0}{m/s}$、そり $M=\qty{20}{kg}$。摩擦力の大きさ $f=\mu'mg$（力-2）。
\phriyu
（1）水平の外力がないので、すべっている時間が長くても運動量保存が使えます（モデル化の札がくずれるのは外力があるときで、ここでは外力がない）。（2）（3）は、それぞれに運動方程式を立て、一定の加速度の式（力-1）で求めます。
\phtateru
\[mv=(m+M)V\qquad \text{荷物：}ma_1=-f,\qquad \text{そり：}Ma_2=f\]
\phtoku
（1）$10\times3.0=30V$ より $V=\qty{1.0}{m/s}$。\\
（2）摩擦力は $f=0.30\times10\times9.8=\qty{29.4}{N}$ なので、
\[a_1=\frac{-29.4}{10}=\qty{-2.94}{m/s^2},\qquad a_2=\frac{29.4}{20}=\qty{1.47}{m/s^2}\]
荷物の速さ $3.0-2.94t$ と、そりの速さ $1.47t$ が等しくなるのは、
\[3.0=4.41t\qquad t\approx\qty{0.68}{s}\]
（3）この間に進む距離は、
\[\text{荷物：}3.0t-\frac{1}{2}\times2.94t^2\approx\qty{1.36}{m},\qquad \text{そり：}\frac{1}{2}\times1.47t^2\approx\qty{0.34}{m}\]
なので、すべった距離は約\qty{1.0}{m}です。
\phtashikame
速さ：$t=0.68$ で荷物 $3.0-2.94\times0.68\approx1.0$、そり $1.47\times0.68\approx1.0$ で、（1）と合います。\\
エネルギー：運動エネルギーは $45-15=\qty{30}{J}$ 減りました。これは「摩擦力×すべった距離」$29.4\times1.02\approx\qty{30}{J}$ と同じです（摩擦が熱に変えたエネルギー。力-4 で学ぶ考え方）。
\kotae （1）\qty{1.0}{m/s}　（2）約\qty{0.68}{s}　（3）約\qty{1.0}{m}
\end{phreidai}
\begin{jibun}
なめらかな氷の上に止まっている質量\qty{30}{kg}のそりに、質量\qty{10}{kg}の荷物が速さ\qty{4.0}{m/s}ですべりこみました。動摩擦係数を0.40として、いっしょに動くときの速さ、そうなるまでの時間、すべった距離を求めましょう。
\bkjans{\qty{1.0}{m/s}、約\qty{0.77}{s}、約\qty{1.5}{m}}
\end{jibun}

%% ======================== まちがえやすいところ ========================
\Needspace{18\baselineskip}\bksection{tsumazuki}{まちがえやすいところをチェック}
\begin{tsumazuki}{速度の向き（符号）を忘れる}
\bkng{\qty{40}{m/s}で来たボールを、\qty{50}{m/s}で打ち返した。運動量の変化は、$0.145\times(50-40)$。}
\bkok{打ち返す向きを＋にすると、はじめの速度は $-40$。だから、変化は $0.145\times\bigl(50-(-40)\bigr)$。}
運動量は向きをもつ量です。向きが反対になると、止める分と、反対に飛ばす分の両方が変化に入ります（例題1）。
\end{tsumazuki}

\begin{tsumazuki}{運動量が保存すれば、運動エネルギーも保存する？}
\bkng{貨車が連結しても運動量が変わらないのだから、運動エネルギーも変わらない。}
\bkok{運動エネルギーが保存するのは $e=1$ のときだけ。合体（$e=0$）では、いちばん多く失われる。}
運動量保存は外力がなければいつも成り立ちますが、運動エネルギーは変形や熱に変わります（④・例題4）。
\end{tsumazuki}

\begin{tsumazuki}{反発係数を速さの比だけで計算する}
\bkng{動いている相手とぶつかったときも、$e=\dfrac{\text{あとの自分の速さ}}{\text{前の自分の速さ}}$。}
\bkok{$e=-\dfrac{v_1'-v_2'}{v_1-v_2}$。2つの物の速度の差（相対速度）で考える。}
$e=\dfrac{v'}{v}$ でよいのは、相手が動かない床やかべのときだけです。
\end{tsumazuki}

%% ======================== 工学ではこう使う ========================
\Needspace{18\baselineskip}\bksection{phgear}{工学ではこう使う}
\begin{phkougaku}[ぶつかっても乗る人を守る車]
自動車の安全の試験（JNCAP）では、車を時速\qty{55}{km}でコンクリートのかべに正面からぶつけて、乗る人の人形が受ける力をはかります。時速\qty{55}{km}は約\qty{15.3}{m/s}です。

車の前の部分が\qty{0.6}{m}つぶれながら、同じ割合で減速して止まるとすると、力-1の $v^2=2ax$ から、減速は $\dfrac{15.3^2}{2\times0.6}\approx\qty{195}{m/s^2}$（$g$ の約20倍）、止まるまでの時間は約\qty{0.08}{s}です。もし\qty{0.1}{m}しかつぶれなければ、減速は約6倍の、$g$ の約120倍になります。運動量の変化は同じでも、\textbf{つぶれる距離と時間を長くして、力を小さくする}のが、衝突安全の基本の考えです。車の前と後ろの部分（クラッシャブルゾーン）はわざとつぶれやすく作り、乗る人のいる部屋はつぶれにくく作ります。この見積もりの数（\qty{0.6}{m}など）は説明のための仮の値で、本物の車では、つぶれ方をコンピュータで計算し（力-16）、試験で確かめます。

\begin{phketa}[衝突と噴き出す速さ]
\phketaline{バットとボールがふれている時間}{約\qty{0.7}{ms}}
\phketaline{自動車の正面衝突の試験の速さ（JNCAP）}{時速\qty{55}{km}（約\qty{15}{m/s}）}
\phketaline{テニスボールとコンクリートの反発係数}{0.73〜0.76}
\phketaline{液体水素エンジン LE-7A の比推力}{\qty{440}{s}（$u\approx\qty{4.3}{km/s}$）}
\phboxhead{答えが桁で合っているか}
ぶつかる時間が「ミリ秒」なら、平均の力は物の重さの何千倍にもなります。ロケットの $\Delta v$ は「噴き出す速さ $u$ の数倍まで」が目安で、それより大きい答えが出たら、質量比か単位を見直そう。
\end{phketa}

\phboxhead{比推力}
ロケットのエンジンの性能は、ふつう\textbf{比推力}（ひすいりょく）で表します。単位は秒です。噴き出す速さ $u$ を、標準の重力加速度（約\qty{9.81}{m/s^2}）でわった値です。LE-7A の比推力は\qty{440}{s}なので、
\[u=440\times9.81\approx\qty{4.3e3}{m/s}\]
です（例題13）。
\end{phkougaku}
\phstation{衝突安全,ロケット工学,流体機械}{力積、反発係数、衝撃荷重、重心、推力、比推力、質量比}

\begin{column}{「ロケットは空気を押して進む」という誤解}
アメリカのロバート・ゴダードは、1920年ごろ、ロケットで月まで届く可能性を論文に書きました。すると、ニューヨーク・タイムズ紙は1920年1月13日の社説で、「押し返す空気のない宇宙で、ロケットが進めるはずがない」という意味のことを書き、ゴダードを批判しました。運動量保存（作用・反作用）を正しく理解していなかったのです。49年後の1969年7月17日、アポロ11号が月へ向かっている最中に、同じ新聞は「ロケットが真空中でも大気中と同じように働くことは、今でははっきり確かめられた。本紙はあやまりを残念に思う」という訂正記事をのせました。ロケットの式そのものは、1903年にロシアのツィオルコフスキーが発表した論文で広く知られるようになりました（それより前の1810年代に、イギリスでも導かれていました）。
\end{column}

%% ======================== 章末問題 ========================
\bkendproblems
\bkproblemgroup{基本}{4}
\bkproblem 質量\qty{0.27}{kg}のバレーボールが速さ\qty{10}{m/s}で飛んできて、反対向きに速さ\qty{15}{m/s}で打ち返されました。ふれていた時間を\qty{0.020}{s}として、平均の力の大きさを求めましょう。
\bkpans{運動量の変化 $0.27\times25=6.75$、$6.75\div0.020$ より約\qty{3.4e2}{N}}
\bkproblem 質量\qty{0.15}{kg}のボールを、速さ\qty{20}{m/s}から止めます。止まるまでの時間が\qty{0.010}{s}のときと、ミットを引いて\qty{0.050}{s}にしたときの、平均の力の大きさを求めましょう。
\bkpans{\qty{300}{N}と\qty{60}{N}}
\bkproblem なめらかな床の上で、質量\qty{2.0}{kg}の台車が速さ\qty{3.0}{m/s}で、止まっている質量\qty{1.0}{kg}の台車にぶつかって連結しました。連結後の速さと、前後の運動エネルギーを求めましょう。
\bkpans{\qty{2.0}{m/s}。\qty{9.0}{J}から\qty{6.0}{J}}
\bkproblem 氷の上に止まっている質量\qty{60}{kg}の人が、質量\qty{3.0}{kg}のボールを前へ、氷に対して速さ\qty{8.0}{m/s}で投げました。人はどちら向きに、どんな速さで動き出しますか。
\bkpans{後ろ向きに\qty{0.40}{m/s}}
\bkproblem ボールを高さ\qty{1.6}{m}から落とすと、\qty{0.90}{m}まではね返りました。反発係数を求めましょう。
\bkpans{$\sqrt{0.90\div1.6}=0.75$}
\bkproblem 口の断面積\qty{0.50}{cm^2}のノズルから、水が速さ\qty{20}{m/s}で出ています。かべに直角に当てて横へ流れるとき、かべを押す力を求めましょう（水の密度\qty{1000}{kg/m^3}）。
\bkpans{$1000\times\num{0.50e-4}\times20^2=\qty{20}{N}$}
\bkproblemgroup{チャレンジ}{4}
\bkproblem 質量\qty{3.0}{kg}の台車Aが速さ\qty{2.0}{m/s}で、止まっている質量\qty{1.0}{kg}の台車Bにまっすぐぶつかりました。反発係数を0.60として、ぶつかったあとの2台の速度と、失われた運動エネルギーを求めましょう。
\bkpans{A：\qty{1.2}{m/s}、B：\qty{2.4}{m/s}。失われた運動エネルギー\qty{0.96}{J}}
\bkproblem 止まっている質量\qty{100}{kg}のボートの上で、質量\qty{50}{kg}の人が前へ\qty{2.4}{m}歩きました。ボートは水に対してどれだけ動きますか（水の抵抗は無視）。
\bkpans{後ろへ\qty{0.80}{m}}
\bkproblem 反発係数0.60のボールを、高さ\qty{1.6}{m}から放しました。弾まなくなるまでの時間を求めましょう。
\bkpans{$t_0=\sqrt{2\times1.6\div9.8}\approx\qty{0.571}{s}$、$T=0.571\times\dfrac{1.6}{0.4}\approx\qty{2.3}{s}$}
\bkproblem 噴き出す速さ\qty{2.0}{km/s}のロケットで、質量比が5のときと25のときの、ふやせる速さを求めましょう（$\ln5\approx1.61$、$\ln25\approx3.22$）。質量比を5倍にすると、速さは何倍になりますか。
\bkpans{約\qty{3.2}{km/s}と約\qty{6.4}{km/s}。2倍（5倍ではない）}
\phnyushi
\bkproblem なめらかな水平の床の上に、質量\qty{3.0}{kg}の長い台車が止まっています。質量\qty{1.0}{kg}の小さな物体が、速さ\qty{4.0}{m/s}で台車の上にすべりこみました。物体と台車のあいだの動摩擦係数を0.50とします。（1）やがて物体と台車がいっしょに動くときの速さ、（2）そうなるまでの時間、（3）物体が台車から落ちないためには、台車の上の長さは少なくとも何mいりますか。
\bkpans{（1）\nobreak\qty{1.0}{m/s}。（2）\nobreak 約\qty{0.61}{s}。（3）\nobreak すべる距離 約\qty{1.2}{m}以上}
\bkproblem なめらかな水平面の上で、質量\qty{1.0}{kg}の物体Bと質量\qty{2.0}{kg}の物体Cが、ばね定数\qty{600}{N/m}の軽いばね（自然の長さ）でつながれて止まっています。質量\qty{1.0}{kg}の物体Aが速さ\qty{6.0}{m/s}でBにまっすぐぶつかりました（AとBの反発係数1、ぶつかる時間はとても短い）。（1）ぶつかった直後のAとBの速さ、（2）ばねがいちばん縮んだときのBとCの速さ、（3）そのときのばねの縮みを求めましょう。
\bkpans{（1）\nobreak A：0、B：\qty{6.0}{m/s}。（2）\nobreak \qty{2.0}{m/s}。（3）\nobreak $\dfrac{1}{2}\times600x^2=18-6=12$ より\qty{0.20}{m}}

\begin{furikaeri}
\begin{checklist}
\item 力積と運動量から、ぶつかるときの平均の力を求められる。\bkfurref{ph:力-3:sec1}
\item 運動量保存と反発係数から、衝突のあとの速さを求められる。\bkfurref{ph:力-3:sec4}
\item 質量が変わるロケットの速さを求められる。\bkfurref{ph:力-3:sec6}
\end{checklist}
\end{furikaeri}
\bkfukidashi{sora}{ロケットは、自分で積んできたガスを押して進むんだね。運動量を数えると、ぶつかるときの細かい力がわからなくても答えが出せるよ。}
